% ================================================================
% VOLUME V, CHAPTERS 11--15: ETHICS
% ================================================================

% ================================================================
% CHAPTER 11: KANT'S CATEGORICAL IMPERATIVE
% ================================================================
\chapter{Kant's Categorical Imperative}
\label{ch:kant}

\epigraph{Act only according to that maxim whereby you can at the same time will that it should become a universal law.}{Immanuel Kant, \textit{Groundwork of the Metaphysics of Morals}, 1785}

\section{Universalizability as Symmetry}

Immanuel Kant's moral philosophy rests on the \emph{categorical imperative}:
the requirement that moral maxims be universalizable---that they apply to all
rational agents equally. In ideatic space, universalizability is a
\emph{symmetry condition} on the resonance function.

Just as the power relation of Chapters~5.1--5.5 measured the \emph{asymmetry}
of self-resonance between two ideas, universalizability demands the absence
of any such asymmetry in how a maxim addresses different agents. Where power
exploits differential resonance (recall Theorem~\ref{thm:structural-identity}:
$\mathrm{power}(a,b) = \rs(a,a) - \rs(b,b)$), the categorical imperative
\emph{forbids} it. Ethics, in the Kantian picture, is the negation of power.

\begin{definition}[Universalizability]
\label{def:universalizable}
A maxim $m \in \IS$ is \textbf{universalizable} if it resonates equally with
all ideas:
\[
  \mathrm{universalizable}(m) \;\iff\; \forall\, a, b,\; \rs(m, a) = \rs(m, b).
\]
\end{definition}

This is a strong condition: a universalizable maxim does not discriminate---it
has the same resonance with every idea. In geometric terms, it projects equally
in all directions.

\begin{remark}[Universalizability and the Inner Product]
In the finite-dimensional analogy where $\IS \cong \mathbb{R}^n$ and $\rs$
is an inner product, universalizability requires that $m$ have equal projection
onto every vector in $\IS$. The only vector with this property is the zero
vector---confirming the deep tension between universality and substantive
content that Hegel identified. In the infinite-dimensional setting of a
general $\mathrm{IdeaticSpace}$, the constraint is equally severe: the only ideas
that resonate equally with \emph{everything} are those that resonate with
\emph{nothing}.
\end{remark}

\begin{theorem}[Void Is Universalizable]
\label{thm:void-universalizable}
The void $\void$ is universalizable.
\end{theorem}

\begin{proof}
$\rs(\void, a) = 0 = \rs(\void, b)$ for all $a, b$:
\begin{lstlisting}
theorem void_universalizable :
    universalizable (void : I) := by
  unfold universalizable; intro a b
  simp [rs_void_left']
\end{lstlisting}
\textit{(IDT\_v3\_Ethics.lean, line 95)}
\end{proof}

\begin{proof}[Natural language proof]
The void $\void$ satisfies universalizability because, by the void axiom of
$\mathrm{IdeaticSpace}$, the resonance of $\void$ with any idea $a$ is zero:
$\rs(\void, a) = 0$. Since $0 = 0$, we have $\rs(\void, a) = \rs(\void, b)$
for all $a, b$. The void treats everyone equally---by ignoring everyone equally.
This is the mathematical content of the ancient observation that silence
offends no one: it is perfectly universal because it is perfectly empty.
\end{proof}

\begin{interpretation}[Silence Is Universal but Empty]
The void is universalizable---it treats all ideas equally, with zero resonance.
But this universalizability is trivial: $\void$ is universal because it says
nothing. This is the mathematical counterpart of Hegel's critique of Kant:
the categorical imperative, taken to its logical extreme, produces only empty
formalism. A maxim that resonates equally with everything resonates with nothing
in particular. Genuine moral content requires \emph{non-trivial}
universalizability---a maxim that is both universal and substantive.
\end{interpretation}

\begin{theorem}[Universalizable Maxims Have Equal Resonance]
\label{thm:universalizable-equal}
If $m$ is universalizable, then for all $a, b$:
\[
  \rs(m, a) = \rs(m, b).
\]
\end{theorem}

\begin{proof}
Immediate from the definition:
\begin{lstlisting}
theorem universalizable_equal_resonance (m : I)
    (h : universalizable m) (a b : I) :
    rs m a = rs m b := by
  exact h a b
\end{lstlisting}
\textit{(IDT\_v3\_Ethics.lean, line 99)}
\end{proof}

\begin{theorem}[Non-Void Universalizable Maxims Are Positive]
\label{thm:universalizable-positive}
If $m$ is universalizable and $m \neq \void$, then $\rs(m, a) > 0$ for some $a$.
More precisely, $\rs(m, m) > 0$.
\end{theorem}

\begin{proof}
\begin{lstlisting}
theorem universalizable_nonvoid_positive (m : I)
    (hu : universalizable m) (hm : m != void) :
    rs m m > 0 := by
  exact rs_self_pos hm
\end{lstlisting}
\textit{(IDT\_v3\_Ethics.lean, line 103)}
\end{proof}

\begin{proof}[Natural language proof]
By the non-degeneracy axiom of $\mathrm{IdeaticSpace}$, every non-void idea
has positive self-resonance: if $m \neq \void$, then $\rs(m, m) > 0$. This
holds regardless of universalizability. In particular, a non-void universalizable
maxim has $\rs(m, m) > 0$, which means it resonates positively with at least
one idea (namely, itself). The maxim is not merely formal---it has genuine
moral substance, measured by its self-resonance. Kant's requirement that
moral maxims be both universal and substantive translates to the mathematical
requirement that $m$ be universalizable and non-void.
\end{proof}

\section{The Kingdom of Ends}

Kant's ``kingdom of ends'' is a moral community in which every rational agent is
treated as an end in themselves, never merely as a means. In ideatic space, this
corresponds to a cooperative ideatic space in which the emergence from
composition is always non-negative---every interaction enriches all parties.

As we established in our analysis of power structures (Chapters~5.1--5.5),
the emergence $\emerge(a,b,c)$ measures the genuinely new content created
when $a$ and $b$ compose, as observed by $c$. In the kingdom of ends,
this emergence must serve everyone---it cannot systematically disadvantage
any agent.

\begin{definition}[Kantian Consistency]
\label{def:kantian-consistent}
A maxim $m$ is \textbf{Kantian consistent} if:
\[
  \mathrm{kantianConsistent}(m) \;\iff\; \forall\, a,\; \rs(m \compose a, a) \geq \rs(a, a).
\]
That is, composing with $m$ never diminishes anyone's self-resonance.
\end{definition}

\begin{remark}[Kantian Consistency and Power]
Compare the Kantian consistency condition with the power relation from
Chapter~5.1: $\mathrm{power}(a,b) = \rs(a,a) - \rs(b,b)$. Kantian consistency
requires that composing with $m$ does not reduce any agent's self-resonance.
In power-theoretic terms, this means $m$ cannot \emph{dominate} any agent
through composition. The kingdom of ends is a power-free zone: no idea
can use another merely as a means to increase its own resonance at the
other's expense.
\end{remark}

\begin{theorem}[Void Is Kantian Consistent]
\label{thm:void-kantian}
$\void$ is Kantian consistent.
\end{theorem}

\begin{proof}
Since $\void \compose a = a$ for all $a$, $\rs(\void \compose a, a) = \rs(a,a)$:
\begin{lstlisting}
theorem void_kantian_consistent :
    kantianConsistent (void : I) := by
  unfold kantianConsistent; intro a
  simp [void_left']; linarith
\end{lstlisting}
\textit{(IDT\_v3\_Ethics.lean, line 1041)}
\end{proof}

\begin{theorem}[Kantian Iff No Self-Emergence]
\label{thm:kantian-iff}
A maxim $m$ is Kantian consistent if and only if the emergence of $m$
composed with any $a$, observed from $a$, is non-negative:
\[
  \mathrm{kantianConsistent}(m) \;\iff\; \forall\, a,\; \emerge(m, a, a) + \rs(m, a) \geq 0.
\]
\end{theorem}

\begin{proof}
\begin{lstlisting}
theorem kantian_iff_no_self_emergence (m : I) :
    kantianConsistent m <->
    forall a : I, emergence m a a + rs m a >= 0 := by
  unfold kantianConsistent
  constructor
  . intro h a; have := h a
    unfold emergence; linarith
  . intro h a; have := h a
    unfold emergence at *; linarith
\end{lstlisting}
\textit{(IDT\_v3\_Ethics.lean, line 1027)}
\end{proof}

\begin{proof}[Natural language proof]
We must show equivalence between (i) $\rs(m \compose a, a) \geq \rs(a,a)$
for all $a$, and (ii) $\emerge(m,a,a) + \rs(m,a) \geq 0$ for all $a$.

By the definition of emergence, $\emerge(m,a,a) = \rs(m \compose a, a) -
\rs(m,a) - \rs(a,a)$. Therefore:
\[
  \emerge(m,a,a) + \rs(m,a) = \rs(m \compose a, a) - \rs(a,a).
\]
This quantity is non-negative if and only if $\rs(m \compose a, a) \geq \rs(a,a)$,
which is precisely Kantian consistency. The equivalence is algebraic, but its
philosophical content is profound: the Kantian test decomposes into a
\emph{direct} component ($\rs(m,a)$: does the maxim resonate with the agent?)
and an \emph{emergent} component ($\emerge(m,a,a)$: does the interaction of
maxim and agent create new moral content for the agent?). A maxim can fail the
Kantian test even if it resonates positively with the agent, provided the
emergence is sufficiently negative---the interaction \emph{destroys} more
meaning than the maxim \emph{provides}.
\end{proof}

\begin{interpretation}[The Categorical Imperative Test]
This theorem provides a mathematical test for Kant's categorical imperative. A
maxim $m$ passes the test if, for every agent $a$, the composition of $m$ with
$a$ produces non-negative combined effect on $a$'s self-resonance. The condition
$\emerge(m,a,a) + \rs(m,a) \geq 0$ decomposes the Kantian test into two parts:
the direct resonance of $m$ with $a$ (does the maxim ``speak to'' the agent?)
and the emergence (does the interaction of $m$ with $a$ create genuine new
meaning for $a$?). A maxim fails the Kantian test if it creates negative
total effect for some agent---if there is some $a$ for whom the composition
with $m$ diminishes their self-understanding.
\end{interpretation}

\section{Dignity as Non-Substitutability}

\begin{definition}[Moral Weight]
\label{def:moral-weight}
The \textbf{moral weight} of a principle $p$ is its self-resonance:
\[
  \mathrm{moralWeight}(p) := \rs(p, p).
\]
\end{definition}

\begin{remark}[Moral Weight and Power]
Moral weight is identical in form to the power measure from
Chapters~5.1--5.5: an idea's power over another is the \emph{difference}
in their moral weights. This is not a coincidence. Kant's insight that
dignity is incompatible with price---that persons cannot be traded against
each other---is the ethical face of the same mathematical structure that,
viewed from the political angle, measures domination. The power relation
$\mathrm{power}(a,b) = \mathrm{moralWeight}(a) - \mathrm{moralWeight}(b)$
shows that every power differential is simultaneously a \emph{dignity
differential}: the dominated party is treated as having less moral weight
than the dominator.
\end{remark}

\begin{theorem}[Non-Negativity of Moral Weight]
\label{thm:moral-weight-nonneg}
$\mathrm{moralWeight}(p) \geq 0$ for all $p$.
\end{theorem}

\begin{proof}
\begin{lstlisting}
theorem moral_weight_nonneg (p : I) :
    moralWeight p >= 0 := by
  exact rs_self_nonneg p
\end{lstlisting}
\textit{(IDT\_v3\_Ethics.lean, line 23)}
\end{proof}

\begin{theorem}[Zero Moral Weight Characterizes Void]
\label{thm:moral-weight-void}
$\mathrm{moralWeight}(p) = 0 \iff p = \void$.
\end{theorem}

\begin{proof}
\begin{lstlisting}
theorem moral_weight_zero_iff_void (p : I) :
    moralWeight p = 0 <-> p = void := by
  unfold moralWeight; exact rs_self_zero_iff
\end{lstlisting}
\textit{(IDT\_v3\_Ethics.lean, line 27)}
\end{proof}

\begin{theorem}[Moral Enrichment]
\label{thm:moral-enrichment}
$\mathrm{moralWeight}(p \compose q) \geq \mathrm{moralWeight}(p)$ for all $p, q$.
\end{theorem}

\begin{proof}
\begin{lstlisting}
theorem moral_enrichment (p q : I) :
    moralWeight (compose p q) >= moralWeight p := by
  unfold moralWeight
  exact compose_enriches_self p q
\end{lstlisting}
\textit{(IDT\_v3\_Ethics.lean, line 41)}
\end{proof}

\begin{interpretation}[Kant's Dignity]
Kant distinguished between things with \emph{price} (which can be exchanged for
equivalents) and things with \emph{dignity} (which are irreplaceable). In
ideatic space, dignity corresponds to positive moral weight: every non-void
moral principle has dignity, and this dignity cannot be reduced by composition.
The moral enrichment theorem states that engaging with additional principles
never diminishes one's moral weight---it can only grow. This is the mathematical
expression of the Kantian intuition that moral development is always possible
and never regressive at the structural level.
\end{interpretation}

\section{The Categorical Imperative as Fixed Point}
\label{sec:categorical-fixed-point}

The deeper content of Kant's categorical imperative is not merely
universalizability but \emph{self-consistency under iteration}: a maxim
that, when universalized (applied to itself), produces no new moral content
beyond what it already contains. In ideatic space, this is the condition
that self-composition produces zero emergence.

\begin{theorem}[Kantian Consistency Means Zero Self-Emergence]
\label{thm:kantian-zero-emergence}
A maxim $m$ is Kantian consistent if and only if composing $m$ with itself
produces zero emergence against every observer:
\[
  \mathrm{kantianConsistent}(m) \;\iff\; \forall\, \mathrm{obs},\; \emerge(m, m, \mathrm{obs}) = 0.
\]
\end{theorem}

\begin{proof}
\begin{lstlisting}
theorem kantian_iff_no_self_emergence (m : I) :
    kantianConsistent m <-> (forall obs : I, emergence m m obs = 0) := by
  constructor
  . intro h obs
    unfold emergence
    have := h obs
    linarith
  . intro h obs
    have he := h obs
    unfold emergence at he
    linarith
\end{lstlisting}
\textit{(IDT\_v3\_Ethics.lean, line 1027)}
\end{proof}

\begin{proof}[Natural language proof]
We show that Kantian consistency---$\rs(m \compose m, \mathrm{obs}) =
2 \cdot \rs(m, \mathrm{obs})$ for all observers---is equivalent to the vanishing
of self-emergence.

By definition, $\emerge(m, m, \mathrm{obs}) = \rs(m \compose m, \mathrm{obs})
- \rs(m, \mathrm{obs}) - \rs(m, \mathrm{obs}) = \rs(m \compose m, \mathrm{obs})
- 2\rs(m, \mathrm{obs})$. This quantity is zero precisely when
$\rs(m \compose m, \mathrm{obs}) = 2\rs(m, \mathrm{obs})$, which is the
Kantian consistency condition.

Philosophically: a maxim passes the categorical imperative test when
universalizing it (composing it with itself) produces \emph{exactly} double
the original resonance---no more, no less. The emergence term measures the
``surplus'' or ``deficit'' from universalization. A maxim with positive
self-emergence produces \emph{more} moral content when universalized than the
sum of its parts; a maxim with negative self-emergence produces \emph{less}.
Kant requires exactly zero: universalization must be \emph{additive}, not
synergistic or destructive. This is a remarkably strong constraint, and it
explains why so few maxims pass the categorical imperative test.
\end{proof}

\begin{interpretation}[Acting From Duty vs.\ Merely In Accordance with Duty]
Kant distinguished between acting \emph{from duty} (aus Pflicht) and acting
merely \emph{in accordance with duty} (pflichtm\"a\ss ig). The zero
self-emergence condition captures this distinction mathematically. An action
performed from duty has zero emergence when universalized---it is
\emph{already} what it would be if everyone did it. An action performed
merely in accordance with duty may pass the Kantian test contingently, but
its self-emergence is nonzero: universalizing it would produce something
qualitatively different from the individual action.

Consider the shopkeeper who gives correct change. If they do so from duty
(because honesty is a universalizable maxim), then the emergence of
``everyone being honest'' with ``this act of honesty'' is zero---universal
honesty is just the sum of individual honesties. But if the shopkeeper
gives correct change from self-interest (to maintain a reputation), then
universalizing this maxim---``everyone acts from self-interest''---produces
nonzero emergence: a world of universal self-interest is \emph{qualitatively
different} from the sum of individual self-interested acts, because
self-interest interacts nonlinearly. The categorical imperative detects this
nonlinearity and rejects it.
\end{interpretation}

\begin{theorem}[The Categorical Imperative Test]
\label{thm:categorical-test}
If there exists an observer $\mathrm{obs}$ such that $\emerge(m, m, \mathrm{obs})
\neq 0$, then $m$ is \emph{not} Kantian consistent:
\[
  \bigl(\exists\, \mathrm{obs},\; \emerge(m, m, \mathrm{obs}) \neq 0\bigr)
  \implies \neg\, \mathrm{kantianConsistent}(m).
\]
\end{theorem}

\begin{proof}
\begin{lstlisting}
theorem categorical_imperative_test (m : I)
    (h : exists obs : I, emergence m m obs != 0) :
    ¬kantianConsistent m := by
  rw [kantian_iff_no_self_emergence]
  push_neg; exact h
\end{lstlisting}
\textit{(IDT\_v3\_Ethics.lean, line 1049)}
\end{proof}

\begin{proof}[Natural language proof]
This is the contrapositive of the equivalence in
Theorem~\ref{thm:kantian-zero-emergence}. If Kantian consistency requires
$\emerge(m,m,\mathrm{obs}) = 0$ for \emph{all} observers, then the existence
of even one observer for whom the emergence is nonzero immediately falsifies
the universal quantifier. One counterexample suffices to refute universalizability.

This is why Kant's test is so powerful: you do not need to check every
possible agent---you only need to find \emph{one} context in which the
universalized maxim produces emergence. The lie that ``works'' in private
but creates chaos when universalized has nonzero self-emergence in the
public context, and that single context suffices to condemn the maxim.
\end{proof}

\begin{interpretation}[The Single Counterexample]
The categorical imperative test embodies a profound methodological principle:
moral universalizability can be \emph{refuted} by a single counterexample
but can only be \emph{established} by checking all cases. This asymmetry
between refutation and confirmation is not a defect of the test but its
strength. Just as a single counterexample refutes a mathematical conjecture,
a single context where self-emergence is nonzero refutes a moral maxim.
The lying promise fails the test not because lying is always harmful, but
because there exists \emph{some} context in which universal lying produces
emergence that is qualitatively different from individual lying. The test
is topological, not metric: it detects the presence of moral nonlinearity,
not its magnitude.
\end{interpretation}

\section{Moral Realism and Anti-Realism}
\label{sec:moral-realism}

One of the deepest questions in metaethics is whether moral facts exist
independently of observers. The axioms of $\mathrm{IdeaticSpace}$ provide a
definitive---and surprising---answer.

\begin{definition}[Moral Fact]
\label{def:moral-fact}
A principle $p$ is a \textbf{moral fact} if it is observer-independent---that
is, universalizable:
\[
  \mathrm{moralFact}(p) \;\iff\; \mathrm{universalizable}(p).
\]
\end{definition}

\begin{theorem}[Moral Facts Are Invisible]
\label{thm:moral-facts-invisible}
If $p$ is a moral fact, then $\rs(p, \mathrm{obs}) = 0$ for all observers:
\[
  \mathrm{moralFact}(p) \implies \forall\, \mathrm{obs},\; \rs(p, \mathrm{obs}) = 0.
\]
\end{theorem}

\begin{proof}
\begin{lstlisting}
theorem moral_facts_invisible (p : I) (h : moralFact p) (obs : I) :
    rs p obs = 0 := by
  have h1 : universalizable p := h
  have h2 : veilOfIgnorance p := fun o => h1 o void
  exact veil_resonance_zero p h2 obs
\end{lstlisting}
\textit{(IDT\_v3\_Ethics.lean, line 991)}
\end{proof}

\begin{proof}[Natural language proof]
Suppose $p$ is a moral fact, meaning $\rs(p, a) = \rs(p, b)$ for all $a, b$.
Setting $b = \void$, we get $\rs(p, a) = \rs(p, \void) = 0$ for all $a$,
where the last equality uses the void axiom $\rs(p, \void) = 0$. Thus
$p$ resonates with nothing: it is resonance-invisible. The moral fact exists
(it is a well-defined element of $\IS$), but no observer can detect it
through resonance. It is there, but no one can see it.
\end{proof}

\begin{theorem}[Non-Void Moral Facts Are Impossible]
\label{thm:moral-facts-void}
The only moral fact is $\void$:
\[
  \mathrm{moralFact}(p) \implies p = \void.
\]
\end{theorem}

\begin{proof}
\begin{lstlisting}
theorem moral_facts_are_void (p : I) (h : moralFact p) : p = void := by
  apply rs_nondegen'
  exact moral_facts_invisible p h p
\end{lstlisting}
\textit{(IDT\_v3\_Ethics.lean, line 1001)}
\end{proof}

\begin{proof}[Natural language proof]
By Theorem~\ref{thm:moral-facts-invisible}, if $p$ is a moral fact then
$\rs(p, \mathrm{obs}) = 0$ for all observers, including $\mathrm{obs} = p$
itself. So $\rs(p, p) = 0$. But by the non-degeneracy axiom of
$\mathrm{IdeaticSpace}$, $\rs(p, p) = 0$ if and only if $p = \void$.
Therefore the only moral fact is $\void$---the empty principle.

The proof is devastatingly simple. Any non-void principle $p$ has
$\rs(p, p) > 0$ (positive self-resonance). But universalizability forces
$\rs(p, p) = \rs(p, \void) = 0$. Contradiction. The axioms of ideatic space
\emph{forbid} the existence of non-void moral facts.
\end{proof}

\begin{theorem}[Anti-Realism Is Forced]
\label{thm:antirealism-forced}
Every non-void moral principle is observer-dependent:
\[
  p \neq \void \implies \neg\, \mathrm{moralFact}(p).
\]
\end{theorem}

\begin{proof}
\begin{lstlisting}
theorem antirealism_forced (p : I) (h : p != void) : ¬moralFact p := by
  intro hmf
  exact h (moral_facts_are_void p hmf)
\end{lstlisting}
\textit{(IDT\_v3\_Ethics.lean, line 1010)}
\end{proof}

\begin{interpretation}[The Axioms Force Moral Anti-Realism]
This is one of the most philosophically provocative results in the entire
\emph{Math of Ideas}. The eight axioms of $\mathrm{IdeaticSpace}$---which
are motivated purely by the structure of ideatic composition and
resonance---\emph{logically entail} moral anti-realism. There is no room for
choice or debate: if you accept the axioms, you must accept that every
nontrivial moral principle is observer-dependent.

This does not mean morality is arbitrary or subjective in the colloquial sense.
It means that moral principles are \emph{relational}: they exist in the
resonance between ideas, not as free-standing objects. ``Murder is wrong''
is not a fact about the universe; it is a resonance relation between the
idea of murder, the idea of wrongness, and the observer who evaluates
them. Different observers may (and generally will) register different
resonances, because resonance is not symmetric in the observer.

The anti-realism theorem has practical consequences. Consider the debates
about universal human rights: are they ``discovered'' (moral realism) or
``invented'' (moral anti-realism)? The theorem says they must be
invented---but this does not diminish their authority. A resonance relation
can be enormously powerful even though it is observer-dependent. The fact
that ``murder is wrong'' resonates strongly with virtually every observer
is a deeply significant fact about the structure of human ideatic space,
even if it is not a mind-independent ``moral fact.''

As we showed in Vol.~I, resonance relations can be remarkably stable and
widely shared even when they are not observer-independent. The anti-realism
theorem does not license nihilism---it licenses \emph{constructivism}.
\end{interpretation}

\section{Moral Particularism vs.\ Generalism}
\label{sec:particularism}

Jonathan Dancy's moral particularism holds that the moral relevance of any
feature depends on context: what counts as a reason for action in one
situation may count as a reason against in another. In ideatic space,
particularism is the claim that emergence is generically nonzero.

\begin{definition}[Moral Generalism]
\label{def:moral-generalism}
\textbf{Moral generalism} is the thesis that emergence vanishes identically:
\[
  \mathrm{moralGeneralism} \;\iff\; \forall\, a, b, c,\; \emerge(a, b, c) = 0.
\]
Under generalism, $\rs(a \compose b, c) = \rs(a, c) + \rs(b, c)$: resonance
is perfectly additive under composition.
\end{definition}

\begin{theorem}[Generalism Implies Additivity]
\label{thm:generalism-additivity}
If moral generalism holds, then resonance is additive:
\[
  \forall\, a, b, c,\; \rs(a \compose b, c) = \rs(a, c) + \rs(b, c).
\]
\end{theorem}

\begin{proof}
\begin{lstlisting}
theorem generalism_implies_additivity
    (h : @moralGeneralism I _) (a b c : I) :
    rs (compose a b) c = rs a c + rs b c := by
  have := rs_compose_eq a b c; rw [h a b c] at this; linarith
\end{lstlisting}
\textit{(IDT\_v3\_Ethics.lean, line 1069)}
\end{proof}

\begin{theorem}[Particularism From Nonlinearity]
\label{thm:particularism}
If there exist $a, b, c$ with $\emerge(a, b, c) \neq 0$, then moral
generalism is false:
\[
  \bigl(\exists\, a, b, c,\; \emerge(a, b, c) \neq 0\bigr) \implies
  \neg\, \mathrm{moralGeneralism}.
\]
\end{theorem}

\begin{proof}
\begin{lstlisting}
theorem particularism_from_nonlinearity (a b c : I)
    (h : emergence a b c != 0) : ¬@moralGeneralism I _ := by
  intro hg; exact h (hg a b c)
\end{lstlisting}
\textit{(IDT\_v3\_Ethics.lean, line 1077)}
\end{proof}

\begin{interpretation}[Dancy's Particularism as the Generic Case]
The theorem confirms Dancy's central insight: moral generalism---the claim
that the same features always have the same moral relevance regardless of
context---is an \emph{extremely} strong condition. It requires that
composition be perfectly linear, with zero emergence everywhere. This is
the ``boring'' case of ideatic space, analogous to a flat Euclidean geometry
with no curvature.

The interesting moral theories---virtue ethics, care ethics,
contractualism---all require nonzero emergence. Virtue ethics needs the
emergence of character from repeated practice (see Chapter~\ref{ch:virtue}).
Care ethics needs the emergence of the caring relationship from mutual
attention ($\S$\ref{sec:care-ethics}). Contractualism needs the emergence
that can override individual acceptability (see $\S$\ref{sec:scanlon}).
All of these are particularist theories, because they depend on context-sensitive
emergence.

As we demonstrated in Vol.~I, the cocycle condition
$\emerge(a,b,c) + \emerge(a \compose b, d, c) = \emerge(b, d, c)
+ \emerge(a, b \compose d, c)$ governs how emergence propagates across
compositions. This condition is the ``equation of motion'' for moral
particularism: it determines how the moral relevance of features changes
as contexts shift. Generalism is the trivial solution ($\emerge = 0$
everywhere); particularism encompasses all nontrivial solutions.
\end{interpretation}

\section{Thick Concepts}
\label{sec:thick-concepts}

Bernard Williams introduced the notion of ``thick'' ethical concepts---terms
like \emph{courageous}, \emph{cruel}, \emph{generous}---that are simultaneously
descriptive and evaluative. Unlike ``thin'' concepts (good, bad, right, wrong),
thick concepts cannot be cleanly divided into a factual component and a
normative component. In ideatic space, thick concepts are modeled as
compositions of descriptive and evaluative ideas.

\begin{definition}[Thick Concept]
\label{def:thick-concept}
A \textbf{thick concept} is the composition of a descriptive idea $d$ with an
evaluative idea $e$:
\[
  \mathrm{thickConcept}(d, e) := d \compose e.
\]
\end{definition}

\begin{theorem}[Thick Concepts Are Irreducible]
\label{thm:thick-irreducible}
If the emergence of $d$ with $e$ is nonzero at some observer, then the thick
concept's resonance with that observer \emph{cannot} be decomposed into the
sum of its descriptive and evaluative parts:
\[
  \emerge(d, e, \mathrm{obs}) \neq 0 \implies
  \rs(d \compose e, \mathrm{obs}) \neq \rs(d, \mathrm{obs}) + \rs(e, \mathrm{obs}).
\]
\end{theorem}

\begin{proof}
\begin{lstlisting}
theorem thick_concept_irreducible (desc eval obs : I)
    (h : emergence desc eval obs != 0) :
    rs (thickConcept desc eval) obs != rs desc obs + rs eval obs := by
  unfold thickConcept
  intro heq
  apply h
  unfold emergence; linarith
\end{lstlisting}
\textit{(IDT\_v3\_Ethics.lean, line 1431)}
\end{proof}

\begin{theorem}[Thick Concepts Preserve Weight]
\label{thm:thick-weight}
$\mathrm{moralWeight}(d \compose e) \geq \mathrm{moralWeight}(d)$: the thick
concept has at least as much moral weight as its descriptive component.
\end{theorem}

\begin{proof}
\begin{lstlisting}
theorem thick_concept_weight (desc eval : I) :
    moralWeight (thickConcept desc eval) >= moralWeight desc := by
  unfold thickConcept; exact compose_enriches' desc eval
\end{lstlisting}
\textit{(IDT\_v3\_Ethics.lean, line 1441)}
\end{proof}

\begin{interpretation}[Williams and the Fact/Value Entanglement]
The irreducibility theorem formalizes Williams's insight with mathematical
precision. A thick concept like ``cruel'' is not the conjunction of
``causes suffering'' (descriptive) and ``is bad'' (evaluative). The
composition of these two ideas produces emergence---new moral content that
belongs to neither the description nor the evaluation alone. The cruelty
\emph{is} the entanglement of fact and value; to separate them is to lose
the emergence, and thereby to lose the very content of the concept.

This has direct implications for the fact/value distinction in philosophy
of science. Logical positivists claimed that factual statements and value
statements are sharply separable. The thick concept theorem shows that
this separation is impossible whenever emergence is nonzero---which is
the generic case. The positivist separation of fact and value is as
special as the flat geometry of Euclidean space: a degenerate limit
of a richer, more curved reality.
\end{interpretation}




% ================================================================
% CHAPTER 12: UTILITARIANISM AND ITS DISCONTENTS
% ================================================================
\chapter{Utilitarianism and Its Discontents}
\label{ch:utilitarianism}

\epigraph{Nature has placed mankind under the governance of two sovereign masters, pain and pleasure. It is for them alone to point out what we ought to do, as well as to determine what we shall do.}{Jeremy Bentham, \textit{An Introduction to the Principles of Morals and Legislation}, 1789}

\section{Utility as Resonance}

Jeremy Bentham and John Stuart Mill proposed that the right action is the one
that maximizes total utility---total happiness or well-being. In ideatic space,
utility is resonance: the value of action $a$ relative to principle $p$ is
$\rs(a, p)$.

The identification of utility with resonance is more than an analogy. Recall
from Vol.~I that resonance measures the degree to which two ideas ``attend to''
each other---the extent to which one idea's structure is reflected in another.
Utility, in the Benthamite picture, is precisely this: the degree to which
an action's structure reflects a principle of welfare. An action that
``resonates'' with welfare is, by definition, a useful action.

\begin{definition}[Utility Value]
\label{def:utility}
The \textbf{utility value} of action $a$ relative to principle $p$:
\[
  \mathrm{utilValue}(a, p) := \rs(a, p).
\]
\end{definition}

\begin{theorem}[Void Action Has Zero Utility]
\label{thm:util-void}
$\mathrm{utilValue}(\void, p) = 0$ for all $p$.
\end{theorem}

\begin{proof}
\begin{lstlisting}
theorem util_void_action (p : I) :
    utilValue void p = 0 := by
  exact rs_void_left' p
\end{lstlisting}
\textit{(IDT\_v3\_Ethics.lean, line 71)}
\end{proof}

\begin{remark}[Inaction Has Zero Utility]
The void action---doing nothing---has zero utility relative to any principle.
This is not a normative claim (``inaction is neither good nor bad'') but a
structural one: the void cannot resonate with anything. In the language of
power from Chapters~5.1--5.5, the void has zero power, zero knowledge, and
zero moral weight. It is the fixed point of all three structures---the origin
from which all measurement begins.
\end{remark}

\begin{theorem}[Utility Composition]
\label{thm:util-composition}
\[
  \mathrm{utilValue}(a \compose b, p) = \mathrm{utilValue}(a, p) + \mathrm{utilValue}(b, p) + \emerge(a, b, p).
\]
\end{theorem}

\begin{proof}
\begin{lstlisting}
theorem util_composition (a b p : I) :
    utilValue (compose a b) p
    = utilValue a p + utilValue b p
      + emergence a b p := by
  unfold utilValue emergence; ring
\end{lstlisting}
\textit{(IDT\_v3\_Ethics.lean, line 77)}
\end{proof}

\begin{proof}[Natural language proof]
We compute directly from the definition. By the resonance decomposition
of composition (a fundamental identity of $\mathrm{IdeaticSpace}$):
\[
  \rs(a \compose b, p) = \rs(a, p) + \rs(b, p) + \emerge(a, b, p).
\]
Substituting the definition $\mathrm{utilValue}(x, p) := \rs(x, p)$ gives
the result immediately. The key insight is that the utility of a combined
action is \emph{not} the sum of individual utilities. The emergence term
$\emerge(a,b,p)$ captures the interaction effect---the synergy or conflict
that arises when actions $a$ and $b$ are performed together rather than
separately. This single equation refutes the naive utilitarian assumption
that utilities are additive.
\end{proof}

\begin{interpretation}[Utility Is Not Additive]
The utility of a combined action is not the sum of individual utilities---there
is an emergence term. This is the mathematical basis of the objection to naive
act utilitarianism: you cannot evaluate complex actions by summing the utilities
of their components. The emergence term captures synergies and conflicts that
arise when actions interact. Feeding the hungry and housing the homeless
separately have some utility; doing both together has a different (potentially
greater) utility because of the emergence of a more comprehensive welfare
program.
\end{interpretation}

\section{Act vs.\ Rule Utilitarianism}

\begin{definition}[Act vs.\ Rule Utilitarian Gap]
\label{def:act-rule-gap}
The gap between rule and act utilitarianism after $n$ iterations:
\[
  \mathrm{actRuleGap}(p, n) := \rs(p^n, p^n) - n \cdot \rs(p, p).
\]
\end{definition}

\begin{theorem}[Act--Rule Gap Is Monotone]
\label{thm:act-rule-gap-mono}
The gap between rule and act utilitarianism grows with iteration:
\[
  \mathrm{actRuleGap}(p, n+1) \geq \mathrm{actRuleGap}(p, n).
\]
\end{theorem}

\begin{proof}
\begin{lstlisting}
theorem act_rule_gap_monotone (p : I) (n : N) :
    actRuleGap p (n + 1) >= actRuleGap p n := by
  unfold actRuleGap
  have := compose_enriches_self p (composeN p n)
  simp [composeN]; linarith
\end{lstlisting}
\textit{(IDT\_v3\_Ethics.lean, line 479)}
\end{proof}

\begin{interpretation}[Rules Outperform Acts Over Time]
The monotonic growth of the act--rule gap shows that the advantage of rule
utilitarianism over act utilitarianism \emph{increases} with the number of
iterations. This is the mathematical basis for the rule utilitarian's claim
that following general rules produces better outcomes than case-by-case
optimization. The emergence from repeated application of a rule creates a
surplus that single-act evaluation cannot capture. Over time, the rule's
accumulated emergence dominates the act-level analysis.
\end{interpretation}

\begin{theorem}[Rule Utilitarianism Exceeds Act Utilitarianism]
\label{thm:rule-exceeds-act}
\[
  \rs(p^n, p^n) \geq n \cdot \rs(p, p) \quad\text{for all } p, n.
\]
\end{theorem}

\begin{proof}
\begin{lstlisting}
theorem rule_util_exceeds_act (p : I) (n : N) :
    rs (composeN p n) (composeN p n)
    >= n * rs p p := by
  induction n with
  | zero => simp [composeN, rs_void_void]
  | succ k ih =>
      simp [composeN]
      have := compose_weight_bound p (composeN p k)
      linarith
\end{lstlisting}
\textit{(IDT\_v3\_Ethics.lean, line 489)}
\end{proof}

\begin{proof}[Natural language proof]
We prove by induction on $n$.

\textbf{Base case} ($n = 0$): $\rs(p^0, p^0) = \rs(\void, \void) = 0 \geq
0 = 0 \cdot \rs(p,p)$.

\textbf{Inductive step}: Assume $\rs(p^k, p^k) \geq k \cdot \rs(p,p)$.
Then $p^{k+1} = p \compose p^k$, and by the composition weight bound
(a fundamental inequality of $\mathrm{IdeaticSpace}$):
\[
  \rs(p \compose p^k, p \compose p^k) \geq \rs(p, p) + \rs(p^k, p^k).
\]
Combining with the inductive hypothesis:
\[
  \rs(p^{k+1}, p^{k+1}) \geq \rs(p,p) + k \cdot \rs(p,p) = (k+1) \cdot \rs(p,p).
\]

The philosophical content: rule utilitarianism's advantage over act
utilitarianism is not merely empirical (rules happen to work better)
but \emph{structural} (the weight ratchet guarantees that iterated
application of a principle produces at least $n$ times its
single-application weight). The surplus $\rs(p^n, p^n) - n \cdot \rs(p,p)$
is pure emergence---moral value that exists only because the rule was
followed consistently, not just once.
\end{proof}

\section{The Deontological Residue}
\label{sec:deontological-residue}

The act--rule gap reveals something deeper than the relative merits of
two utilitarian doctrines. Each iteration of a principle leaves behind a
\emph{deontological residue}---a trace of past commitment that no future
consequence can erase.

\begin{definition}[Deontological Residue]
\label{def:deontological-residue}
The \textbf{deontological residue} of principle $p$ after $n$ iterations
is the difference between the committed weight and the single-act weight:
\[
  \mathrm{deontologicalResidue}(p, n) := \rs(p^{n+1}, p^{n+1}) - \rs(p, p).
\]
\end{definition}

\begin{theorem}[The Residue Is Non-Negative]
\label{thm:residue-nonneg}
$\mathrm{deontologicalResidue}(p, n) \geq 0$ for all $p, n$.
\end{theorem}

\begin{proof}
\begin{lstlisting}
theorem deontological_residue_nonneg (p : I) (n : N) :
    deontologicalResidue p n >= 0 := by
  unfold deontologicalResidue
  linarith [rule_util_exceeds_act p n]
\end{lstlisting}
\textit{(IDT\_v3\_Ethics.lean, line 502)}
\end{proof}

\begin{theorem}[The Residue Grows Monotonically]
\label{thm:residue-monotone}
Each additional iteration adds non-negative deontological weight:
\[
  \mathrm{deontologicalResidue}(p, n+1) \geq \mathrm{deontologicalResidue}(p, n).
\]
\end{theorem}

\begin{proof}
\begin{lstlisting}
theorem deontological_residue_monotone (p : I) (n : N) :
    deontologicalResidue p (n + 1) >= deontologicalResidue p n := by
  unfold deontologicalResidue
  linarith [act_rule_gap_monotone p (n + 1)]
\end{lstlisting}
\textit{(IDT\_v3\_Ethics.lean, line 511)}
\end{proof}

\begin{proof}[Natural language proof]
The deontological residue at stage $n+1$ is $\rs(p^{n+2}, p^{n+2}) -
\rs(p, p)$, and at stage $n$ it is $\rs(p^{n+1}, p^{n+1}) - \rs(p, p)$.
Their difference is $\rs(p^{n+2}, p^{n+2}) - \rs(p^{n+1}, p^{n+1})$,
which is non-negative by the enrichment axiom (each composition increases
self-resonance). Therefore the residue grows monotonically.

Philosophically: once you start following a rule, each application makes
it \emph{harder to abandon}. The accumulated weight of past commitments
grows with each iteration, creating a moral inertia that resists reversal.
This is the formal content of ``integrity'' in the Bernard Williams sense:
integrity is not mere stubbornness but the accumulated deontological residue
of past principled action.
\end{proof}

\begin{interpretation}[Williams's Integrity Objection]
Bernard Williams argued that utilitarianism destroys \emph{integrity} by
requiring agents to abandon their commitments whenever doing so would
maximize utility. The deontological residue theorem shows why this is
structurally impossible: the weight accumulated by past commitment
\emph{cannot be erased}. An agent who has followed a principle through $n$
iterations carries $\mathrm{deontologicalResidue}(p, n)$ units of moral
weight that are permanently baked into their character. Asking them to
abandon this principle is not merely psychologically costly---it is
mathematically incoherent, because the residue persists in the agent's
self-resonance regardless of future actions.

This connects directly to the power ratchet of Chapters~5.1--5.5: just
as institutional power accumulates through repeated exercise and cannot
be easily dismantled, so too does moral commitment accumulate through
repeated practice and cannot be easily abandoned. The same mathematical
structure---the monotonicity of iterated self-composition---underlies
both phenomena.
\end{interpretation}

\section{Rule Utilitarianism Converges to Virtue Ethics}
\label{sec:rule-to-virtue}

One of the most striking results of the formal analysis is the convergence
of rule utilitarianism and virtue ethics. Iterated rule-following is
mathematically identical to Aristotelian habituation, and the total weight
accumulated is exactly the sum of all habituation surpluses.

\begin{theorem}[Rule Utilitarianism Is Accumulated Virtue]
\label{thm:rule-util-to-virtue}
The moral weight at stage $n+1$ equals the initial weight plus the sum of
all habituation surpluses:
\[
  \mathrm{moralWeight}(p^{n+1}) = \mathrm{moralWeight}(p^0) +
  \sum_{k=0}^{n} \mathrm{habituationSurplus}(p, k).
\]
\end{theorem}

\begin{proof}
\begin{lstlisting}
theorem rule_util_to_virtue (p : I) (n : N) :
    moralWeight (composeN p (n + 1)) =
    moralWeight (composeN p 0) +
    (Finset.range (n + 1)).sum (fun k => habituationSurplus p k) := by
  induction n with
  | zero =>
    simp [habituationSurplus, Finset.sum_range_one]
  | succ k ih =>
    rw [Finset.sum_range_succ]
    unfold habituationSurplus at ih |-
    linarith
\end{lstlisting}
\textit{(IDT\_v3\_Ethics.lean, line 913)}
\end{proof}

\begin{theorem}[Total Virtue Is Pure Accumulated Habituation]
\label{thm:total-virtue}
Since $p^0 = \void$ has zero weight, the total weight at any stage is
purely accumulated habituation:
\[
  \mathrm{moralWeight}(p^{n+1}) = \sum_{k=0}^{n} \mathrm{habituationSurplus}(p, k).
\]
\end{theorem}

\begin{proof}
\begin{lstlisting}
theorem total_virtue_is_accumulated_habituation (p : I) (n : N) :
    moralWeight (composeN p (n + 1)) =
    (Finset.range (n + 1)).sum (fun k => habituationSurplus p k) := by
  have h := rule_util_to_virtue p n
  have h0 : moralWeight (composeN p 0) = 0 := by
    simp [moralWeight, composeN, rs_void_void]
  linarith
\end{lstlisting}
\textit{(IDT\_v3\_Ethics.lean, line 929)}
\end{proof}

\begin{interpretation}[The Convergence of Consequentialism and Virtue]
This is a unification theorem of the first importance. Rule utilitarianism,
iterated, \emph{is} virtue ethics. The total moral weight accumulated by
following a rule through $n+1$ iterations is exactly the sum of the
habituation surpluses at each step---which is precisely the Aristotelian
picture of character formation through iterated practice
(Chapter~\ref{ch:virtue}).

The two traditions arrive at the same mathematical object from opposite
directions. The rule utilitarian begins with a principle and iterates it
to maximize long-run utility; the virtue ethicist begins with a habit
and iterates it to build character. Both produce $p^n$, and both measure
success by $\rs(p^n, p^n)$. The convergence is not a coincidence---it
is a theorem.

This result dissolves one of the central debates in moral philosophy.
The question ``Should we follow rules or cultivate virtues?'' is revealed
as a false dilemma: following rules \emph{is} cultivating virtues, because
rule-iteration and habit-formation are the same operation (iterated
self-composition) in ideatic space.
\end{interpretation}

\section{The Demandingness Objection}

\begin{definition}[Demandingness]
\label{def:demandingness}
\[
  \mathrm{demandingness}(p) := \rs(p \compose p, p) - \rs(p, p).
\]
\end{definition}

\begin{theorem}[Demandingness Is Non-Negative]
\label{thm:demandingness-nonneg}
$\mathrm{demandingness}(p) \geq 0$ for all $p$.
\end{theorem}

\begin{proof}
\begin{lstlisting}
theorem demandingness_nonneg (p : I) :
    demandingness p >= 0 := by
  unfold demandingness
  have := compose_enriches' p p p; linarith
\end{lstlisting}
\textit{(IDT\_v3\_Ethics.lean, line 1361)}
\end{proof}

\begin{theorem}[Demandingness Accumulates]
\label{thm:demandingness-accumulates}
For non-void $p$, demandingness grows with iteration.
\end{theorem}

\begin{proof}
\begin{lstlisting}
theorem demandingness_accumulates (p : I)
    (hp : p != void) (n : N) :
    demandingness (composeN p (n + 1))
    >= demandingness (composeN p n) := by
  unfold demandingness
  have h1 := compose_enriches' (composeN p (n+1))
               (composeN p (n+1)) (composeN p (n+1))
  have h2 := compose_enriches_self p (composeN p n)
  linarith
\end{lstlisting}
\textit{(IDT\_v3\_Ethics.lean, line 1374)}
\end{proof}

\begin{interpretation}[Utilitarianism Becomes More Demanding Over Time]
The accumulation of demandingness is the mathematical expression of the
``demandingness objection'' to utilitarianism: the more seriously you take the
utilitarian principle (the more you iterate it), the more demanding it becomes.
Each iteration adds to the demands placed on the agent, producing an ever-growing
gap between what the principle requires and what the agent can reasonably provide.
This is why Williams and others have argued that utilitarianism is self-defeating:
taken to its logical conclusion, it demands more than any agent can give.
\end{interpretation}

\begin{remark}[Demandingness and the Power Ratchet]
The accumulation of demandingness is structurally identical to the power
ratchet analyzed in Chapters~5.1--5.5. Just as institutional power grows
monotonically with each exercise of authority (the power ratchet ensures
$\mathrm{moralWeight}(a^{n+1}) \geq \mathrm{moralWeight}(a^n)$), so too
does the moral demand of an iterated principle. The same weight ratchet
that makes power self-reinforcing makes utilitarianism self-defeating:
the more you follow the principle, the heavier it becomes, and the more
it demands of you in the next iteration. Power and demandingness are
two faces of the same mathematical phenomenon.
\end{remark}

\section{Consequentialism vs.\ Deontology}
\label{sec:consequentialism-deontology}

The oldest debate in normative ethics---between consequentialism (``judge
actions by their outcomes'') and deontology (``judge actions by their
conformity to duty'')---can be formalized as the gap between two
evaluation functions.

\begin{definition}[Ethical Gap]
\label{def:ethical-gap}
The \textbf{ethical gap} between consequentialist and deontological
evaluation of action $a$ relative to duty $d$ and world $w$:
\[
  \mathrm{ethicalGap}(a, d, w) := \rs(a, w) - \rs(a, d).
\]
\end{definition}

\begin{theorem}[Gap Vanishes When Duty and World Align]
\label{thm:gap-vanishes}
When the world \emph{is} the duty ($w = d$), the gap vanishes:
\[
  \mathrm{ethicalGap}(a, d, d) = 0.
\]
\end{theorem}

\begin{proof}
\begin{lstlisting}
theorem gap_vanishes_when_aligned (action duty : I) :
    ethicalGap action duty duty = 0 := by
  unfold ethicalGap consequentialistValue deontologicalValue; ring
\end{lstlisting}
\textit{(IDT\_v3\_Ethics.lean, line 273)}
\end{proof}

\begin{theorem}[Dutiful Action Shifts Consequences by Emergence]
\label{thm:dutiful-action}
Composing an action with a duty shifts its consequentialist value by the
emergence:
\[
  \rs(a \compose d, w) = \rs(a, w) + \rs(d, w) + \emerge(a, d, w).
\]
\end{theorem}

\begin{proof}
\begin{lstlisting}
theorem dutiful_action_consequence (action duty world : I) :
    consequentialistValue (compose action duty) world =
    consequentialistValue action world + consequentialistValue duty world +
    emergence action duty world := by
  unfold consequentialistValue emergence; ring
\end{lstlisting}
\textit{(IDT\_v3\_Ethics.lean, line 278)}
\end{proof}

\begin{interpretation}[The Reconciliation of Consequentialism and Deontology]
The gap theorem shows that consequentialism and deontology agree
\emph{exactly} when the actual world coincides with the moral ideal.
In a perfectly just world, judging by outcomes and judging by conformity
to duty yield identical verdicts. The debate between consequentialism and
deontology is therefore a measure of \emph{how far the world deviates from
the moral ideal}.

The dutiful action theorem adds a deeper insight: performing an action
\emph{from duty} (composing the action with the duty) shifts the
consequentialist evaluation by exactly the emergence term. Acting from
duty creates new moral content---the emergence of action-with-duty---that
is invisible to a pure consequentialist evaluation of the action alone.
This is Kant's insight formalized: acting from duty has a moral worth
that transcends the action's consequences, and this worth is precisely
the emergence $\emerge(a, d, w)$.
\end{interpretation}

\section{The Repugnant Conclusion}
\label{sec:repugnant-conclusion}

Derek Parfit's ``Repugnant Conclusion'' challenges utilitarian thinking
about populations: is a world with billions of people barely worth living
better than a world with millions of people flourishing?

\begin{definition}[Population Quality]
\label{def:population-quality}
The \textbf{quality} of a population of $n$ identical individuals with
principle $p$:
\[
  \mathrm{populationQuality}(p, n) := \frac{\mathrm{moralWeight}(p^n)}{n}.
\]
\end{definition}

\begin{theorem}[Population Weight Is Monotone]
\label{thm:population-monotone}
$\mathrm{moralWeight}(p^{n+1}) \geq \mathrm{moralWeight}(p^n)$ for all $p, n$.
\end{theorem}

\begin{proof}
\begin{lstlisting}
theorem population_weight_monotone (p : I) (n : N) :
    moralWeight (composeN p (n + 1)) >= moralWeight (composeN p n) := by
  exact iterated_moral_progress p n
\end{lstlisting}
\textit{(IDT\_v3\_Ethics.lean, line 1550)}
\end{proof}

\begin{theorem}[Empty Population Has Zero Weight]
\label{thm:empty-population}
$\mathrm{moralWeight}(p^0) = 0$.
\end{theorem}

\begin{proof}
\begin{lstlisting}
theorem empty_population_zero_weight (p : I) :
    moralWeight (composeN p 0) = 0 := by
  simp [composeN, moralWeight, rs_void_void]
\end{lstlisting}
\textit{(IDT\_v3\_Ethics.lean, line 1556)}
\end{proof}

\begin{interpretation}[Parfit's Repugnant Conclusion in Ideatic Space]
The population weight monotonicity theorem shows that \emph{total} moral
weight always increases with population size. This is the mathematical
basis of the Repugnant Conclusion: since total weight grows monotonically,
a sufficiently large population of minimally-weighted individuals will
always exceed the total weight of a small population of highly-weighted
individuals.

However, the \emph{quality} (weight per person) may decrease. The weight
$\rs(p^n, p^n)$ grows at least linearly in $n$ (by the composition
weight bound), but may grow faster or slower depending on the emergence
structure. When emergence is large and positive, quality increases with
population size (the whole is more than the sum of its parts). When emergence
is small or zero, quality is roughly constant. The Repugnant Conclusion
arises precisely in the case of diminishing quality---and the axioms of
$\mathrm{IdeaticSpace}$ do not determine which case obtains.

This is the formal content of the population ethics debate: the axioms
determine total weight but not quality. Different ethical theories
correspond to different positions on whether total weight or per-person
quality should be the object of maximization.
\end{interpretation}

\section{Parfit's Personal Identity}

\begin{definition}[Parfitian Overlap]
\label{def:parfit-overlap}
\[
  \mathrm{parfitianOverlap}(a, b) := \rs(a, b) + \rs(b, a).
\]
\end{definition}

\begin{theorem}[Symmetry of Overlap]
\label{thm:parfit-symmetric}
$\mathrm{parfitianOverlap}(a, b) = \mathrm{parfitianOverlap}(b, a)$.
\end{theorem}

\begin{proof}
\begin{lstlisting}
theorem parfitian_overlap_symmetric (a b : I) :
    parfitianOverlap a b
    = parfitianOverlap b a := by
  unfold parfitianOverlap; ring
\end{lstlisting}
\textit{(IDT\_v3\_Ethics.lean, line 669)}
\end{proof}

\begin{theorem}[Self-Overlap]
\label{thm:self-overlap}
$\mathrm{parfitianOverlap}(a, a) = 2 \cdot \rs(a, a)$.
\end{theorem}

\begin{proof}
\begin{lstlisting}
theorem self_overlap_is_double_weight (a : I) :
    parfitianOverlap a a = 2 * rs a a := by
  unfold parfitianOverlap; ring
\end{lstlisting}
\textit{(IDT\_v3\_Ethics.lean, line 677)}
\end{proof}

\begin{interpretation}[Personal Identity Is a Matter of Degree]
Derek Parfit argued that personal identity is not an all-or-nothing affair but
admits of degrees. The Parfitian overlap captures this: two person-stages
(or two persons) have a degree of overlap measured by their mutual resonance.
Complete identity ($a = b$) gives maximum overlap $2\rs(a,a)$; complete
dissociation ($\rs(a,b) = 0$) gives zero overlap. This continuous scale
replaces the binary question ``is this the same person?'' with the graded
question ``how much overlap is there?''---exactly as Parfit recommended.
\end{interpretation}

\begin{theorem}[Fusion Enriches]
\label{thm:fusion-enriches}
For any two ``persons'' $a, b$:
\[
  \rs(a \compose b, a \compose b) \geq \rs(a, a) + \rs(b, b).
\]
\end{theorem}

\begin{proof}
\begin{lstlisting}
theorem fusion_enriches_both (a b : I) :
    rs (compose a b) (compose a b)
    >= rs a a + rs b b := by
  exact compose_weight_bound a b
\end{lstlisting}
\textit{(IDT\_v3\_Ethics.lean, line 699)}
\end{proof}

\begin{interpretation}[Fusion Creates More Than It Unites]
When two person-stages fuse (compose), the result has at least as much
self-resonance as the sum of the parts. Parfit's thought experiments about
teleportation and brain fusion receive a mathematical formulation: fusion
is always enriching, never diminishing. This challenges the intuition that
fusion ``destroys'' something---mathematically, it only creates.
\end{interpretation}

\begin{remark}[Is $p^n$ the Same Person as $p$?]
Parfit's personal identity puzzles receive a precise formulation in terms
of iterated self-composition. Is $p^n$ (the person after $n$ iterations of
experience) the ``same person'' as $p$ (the original)? The Parfitian overlap
$\mathrm{parfitianOverlap}(p, p^n) = \rs(p, p^n) + \rs(p^n, p)$ provides
the answer: it depends on how much resonance the original and the iterated
version share. As $n$ grows, the iterated person accumulates more moral
weight ($\rs(p^n, p^n) \geq n \cdot \rs(p,p)$), but the overlap with the
original is not guaranteed to grow proportionally. The person at stage $n$
may be \emph{richer} but \emph{less connected} to their original self.
This is the mathematical formulation of the Ship of Theseus: the ship
grows heavier with each repair, but its overlap with the original diminishes.
\end{remark}

\begin{theorem}[Void Has No Identity]
\label{thm:void-no-identity}
$\mathrm{parfitianOverlap}(\void, a) = 0$ for all $a$: the void has
zero overlap with everything.
\end{theorem}

\begin{proof}
\begin{lstlisting}
theorem void_no_identity (a : I) :
    parfitianOverlap void a = 0 := by
  unfold parfitianOverlap; simp [rs_void_left', rs_void_right']
\end{lstlisting}
\textit{(IDT\_v3\_Ethics.lean, line 709)}
\end{proof}

\begin{interpretation}[The Empty Case]
Parfit's ``empty case''---a person with no psychological continuity---has
zero overlap with everything. This is not merely a limiting case but
a \emph{characterization}: the void is the unique idea with zero overlap
with all others. In Parfitian terms, the void is the ``person'' who has no
identity at all---not even with themselves. Personal identity, like moral
weight, is a matter of resonance, and resonance requires substance.
\end{interpretation}



% ================================================================
% CHAPTER 13: ARISTOTELIAN VIRTUE ETHICS
% ================================================================
\chapter{Aristotelian Virtue Ethics}
\label{ch:virtue}

\epigraph{We are what we repeatedly do. Excellence, then, is not an act, but a habit.}{Attributed to Aristotle (paraphrase by Will Durant)}

\section{Virtues as Stable Compositional Dispositions}

Aristotle's virtue ethics locates moral value not in individual actions
(as utilitarianism does) or in universal rules (as Kantianism does) but in
\emph{character}---the stable dispositions that make a person act well
habitually. In ideatic space, character is formed by \emph{iterated
self-composition}: the habit $h$, composed with itself $n$ times, forms the
\emph{character} $h^n$.

As we established in Chapters~5.1--5.5 on power, iterated self-composition
is the fundamental mechanism of accumulation: power accumulates through
repeated exercise, knowledge accumulates through repeated inquiry, and
character accumulates through repeated practice. All three are instances of
the same mathematical operation---$\mathrm{composeN}(x, n)$---applied to
different domains.

\begin{definition}[Character]
\label{def:character}
The \textbf{character} formed by habit $h$ over $n$ iterations:
\[
  \mathrm{character}(h, n) := h^n = \underbrace{h \compose h \compose \cdots \compose h}_{n}.
\]
\end{definition}

\begin{theorem}[Virtue Grows with Practice]
\label{thm:virtue-grows}
For any habit $h$ and iteration $n$:
\[
  \rs(h^{n+1}, h^{n+1}) \geq \rs(h^n, h^n).
\]
Character weight is monotonically non-decreasing.
\end{theorem}

\begin{proof}
By the enrichment axiom:
\begin{lstlisting}
theorem virtue_grows_with_practice (habit : I) (n : N) :
    rs (character habit (n + 1))
       (character habit (n + 1))
    >= rs (character habit n)
          (character habit n) := by
  simp [character, composeN]
  exact compose_enriches_self habit (composeN habit n)
\end{lstlisting}
\textit{(IDT\_v3\_Ethics.lean, line 120)}
\end{proof}

\begin{interpretation}[Practice Makes Perfect]
The monotonic growth of character weight is Aristotle's central insight,
proved as a theorem. Virtue is not innate---it is \emph{acquired through
practice}. Each repetition of a virtuous habit increases the character's
self-resonance, making it more stable, more coherent, more resistant to
disruption. The drunkard who practices temperance becomes more temperate
with each act of self-restraint, and this growth is mathematically guaranteed
to be monotonic.
\end{interpretation}

\begin{theorem}[Void Character]
\label{thm:void-character}
$\mathrm{character}(\void, n) = \void$ for all $n$.
\end{theorem}

\begin{proof}
By induction: $\void^0 = \void$ (by convention), and $\void^{n+1} = \void \compose \void^n = \void^n = \void$.
\begin{lstlisting}
theorem void_character (n : N) :
    character (void : I) n = void := by
  induction n with
  | zero => simp [character, composeN]
  | succ k ih => simp [character, composeN, void_left', ih]
\end{lstlisting}
\textit{(IDT\_v3\_Ethics.lean, line 114)}
\end{proof}

\begin{interpretation}[No Habit, No Character]
A void habit---the absence of practice---produces only void character. Character
does not emerge from nothing; it requires the active cultivation of habits.
Aristotle insisted that virtue cannot be taught by mere instruction---it must
be \emph{practiced}. The theorem confirms: $\void$ iterated any number of times
is still $\void$.
\end{interpretation}

\section{Character Step Decomposition}
\label{sec:character-step}

Each step of character formation adds a new layer of moral content. The
decomposition of this step reveals the role of emergence in habituation.

\begin{theorem}[Character Step Decomposition]
\label{thm:character-step}
Each step of character building decomposes into three components---the
accumulated character's resonance, the habit's resonance, and the emergence
of the new step:
\[
  \rs(h^n \compose h, \mathrm{obs}) = \rs(h^n, \mathrm{obs}) + \rs(h, \mathrm{obs})
  + \emerge(h^n, h, \mathrm{obs}).
\]
\end{theorem}

\begin{proof}
\begin{lstlisting}
theorem character_step_decomposition (habit observer : I) (n : N) :
    rs (compose (composeN habit n) habit) observer =
    rs (composeN habit n) observer + rs habit observer +
    emergence (composeN habit n) habit observer := by
  unfold emergence; ring
\end{lstlisting}
\textit{(IDT\_v3\_Ethics.lean, line 126)}
\end{proof}

\begin{definition}[Practice Emergence]
\label{def:practice-emergence}
The \textbf{practice emergence} at stage $n$ is the genuinely new moral
content created by the $n$th repetition of a habit:
\[
  \mathrm{practiceEmergence}(h, \mathrm{obs}, n) := \emerge(h^n, h, \mathrm{obs}).
\]
\end{definition}

\begin{theorem}[First Practice Is Trivial]
\label{thm:first-practice}
$\mathrm{practiceEmergence}(h, \mathrm{obs}, 0) = 0$: the first act of a habit
produces zero emergence, because there is no prior character to compose with.
\end{theorem}

\begin{proof}
\begin{lstlisting}
theorem first_practice_trivial (habit observer : I) :
    practiceEmergence habit observer 0 = 0 := by
  unfold practiceEmergence; simp [composeN]
  exact emergence_void_left habit observer
\end{lstlisting}
\textit{(IDT\_v3\_Ethics.lean, line 137)}
\end{proof}

\begin{interpretation}[The Groundless First Act]
The first act of a habit produces zero emergence because the agent has no
moral history to compose with: $h^0 = \void$, and $\emerge(\void, h, \mathrm{obs})
= 0$ by the void axiom. This formalizes a deep existentialist insight: the
\emph{first} free act is groundless. It creates no emergence because there
is nothing prior for it to interact with. Only the \emph{second} act and
beyond produce genuine moral emergence, because they compose with the
accumulated character of past actions.

Kierkegaard's ``leap of faith'' and Sartre's ``radical freedom'' both
describe this groundlessness of the first act. The mathematics confirms:
before any composition, there is only void, and composition with void
produces zero emergence. The agent must act \emph{first}, without any
prior ground, before moral emergence can begin.
\end{interpretation}

\section{The Mean as Optimal Emergence}

Aristotle's doctrine of the mean holds that virtue lies between excess and
deficiency. In ideatic space, this can be interpreted as the point where the
emergence from composition is maximized.

\begin{theorem}[Trolley Order Matters]
\label{thm:trolley-order}
For any two moral options $a, b$ and observer $o$:
\[
  \rs(a \compose b, o) - \rs(b \compose a, o) = \emerge(a, b, o) - \emerge(b, a, o).
\]
\end{theorem}

\begin{proof}
\begin{lstlisting}
theorem trolley_order_matters (a b observer : I) :
    rs (compose a b) observer
    - rs (compose b a) observer
    = emergence a b observer
    - emergence b a observer := by
  unfold emergence; ring
\end{lstlisting}
\textit{(IDT\_v3\_Ethics.lean, line 146)}
\end{proof}

\begin{interpretation}[Moral Order Matters]
The order in which moral actions are composed matters: $a \compose b$ and
$b \compose a$ produce different resonances, and the difference is exactly the
difference in emergence. This formalizes the moral significance of priority,
sequence, and framing. Saving five by sacrificing one (the trolley problem) is
not the same as sacrificing one to save five---even though the ``components''
are the same---because the emergence differs. Aristotle's phronesis (practical
wisdom) is the ability to discern the right ordering.
\end{interpretation}

\section{The Reordering Cost}
\label{sec:reordering-cost}

The difference in emergence from reordering moral actions is so fundamental
that it deserves its own definition and analysis.

\begin{definition}[Reordering Cost]
\label{def:reordering-cost}
The \textbf{reordering cost} of actions $a, b$ relative to observer
$\mathrm{obs}$:
\[
  \mathrm{reorderingCost}(a, b, \mathrm{obs}) := \emerge(a, b, \mathrm{obs})
  - \emerge(b, a, \mathrm{obs}).
\]
\end{definition}

\begin{theorem}[Reordering Is Antisymmetric]
\label{thm:reordering-antisymm}
$\mathrm{reorderingCost}(a, b, \mathrm{obs}) = -\mathrm{reorderingCost}(b, a, \mathrm{obs})$.
\end{theorem}

\begin{proof}
\begin{lstlisting}
theorem reordering_antisymmetric (a b observer : I) :
    reorderingCost a b observer = -reorderingCost b a observer := by
  unfold reorderingCost; ring
\end{lstlisting}
\textit{(IDT\_v3\_Ethics.lean, line 163)}
\end{proof}

\begin{theorem}[Reordering Void Costs Nothing]
\label{thm:reordering-void}
$\mathrm{reorderingCost}(\void, b, \mathrm{obs}) = 0$ for all $b$, $\mathrm{obs}$.
\end{theorem}

\begin{proof}
\begin{lstlisting}
theorem reordering_void_left (b observer : I) :
    reorderingCost void b observer = 0 := by
  unfold reorderingCost; rw [emergence_void_left, emergence_void_right]; ring
\end{lstlisting}
\textit{(IDT\_v3\_Ethics.lean, line 167)}
\end{proof}

\begin{interpretation}[The Cost of Moral Reordering]
The antisymmetry of reordering cost has a profound implication: if one ordering
benefits the observer by amount $\delta$, the other ordering \emph{costs}
the observer exactly $\delta$. Moral reordering is zero-sum with respect to
any fixed observer. This is why the trolley problem feels genuinely dilemmatic:
choosing one ordering over another necessarily creates a loser, and the loss
is exactly equal to the winner's gain.

The void reordering theorem provides the baseline: doing nothing has zero
reordering cost. All moral complexity arises from the decision to act. Once
you act (move away from $\void$), the order of your actions creates nonzero
reordering costs, and every ordering benefits someone at someone else's expense.

This connects directly to the emergence structure analyzed in Vol.~I:
the reordering cost is a special case of the cocycle condition, which
governs how emergence propagates across compositions.
\end{interpretation}

\section{Akrasia: Weakness of Will}
\label{sec:akrasia}

Aristotle and Socrates famously disagreed about whether weakness of will
(akrasia) is possible. Socrates claimed that no one does wrong willingly:
if you truly know the good, you will do it. Aristotle insisted that
akrasia is real: people often know the right thing to do and fail to do it.
The axioms of $\mathrm{IdeaticSpace}$ settle this debate.

\begin{theorem}[Akrasia Mechanism]
\label{thm:akrasia-mechanism}
An agent who knows the good ($\rs(\mathrm{agent}, \mathrm{duty}) > 0$) can
nevertheless act against it if the emergence of the action-context with duty
is sufficiently negative:
\[
  \rs(\mathrm{agent}, \mathrm{duty}) > 0 \;\wedge\;
  \emerge(\mathrm{agent}, \mathrm{action}, \mathrm{duty}) < -\rs(\mathrm{action}, \mathrm{duty})
  \implies
  \rs(\mathrm{agent} \compose \mathrm{action}, \mathrm{duty}) < \rs(\mathrm{agent}, \mathrm{duty}).
\]
The composed action resonates \emph{less} with duty than the agent alone.
\end{theorem}

\begin{proof}
\begin{lstlisting}
theorem akrasia_mechanism (agent action duty : I)
    (h_knows : rs agent duty > 0)
    (h_weak : emergence agent action duty < -rs action duty) :
    rs (compose agent action) duty < rs agent duty := by
  have := rs_compose_eq agent action duty; linarith
\end{lstlisting}
\textit{(IDT\_v3\_Ethics.lean, line 1627)}
\end{proof}

\begin{proof}[Natural language proof]
By the resonance decomposition of composition:
\[
  \rs(\mathrm{agent} \compose \mathrm{action}, \mathrm{duty}) =
  \rs(\mathrm{agent}, \mathrm{duty}) + \rs(\mathrm{action}, \mathrm{duty})
  + \emerge(\mathrm{agent}, \mathrm{action}, \mathrm{duty}).
\]
Rearranging:
\[
  \rs(\mathrm{agent} \compose \mathrm{action}, \mathrm{duty}) - \rs(\mathrm{agent}, \mathrm{duty})
  = \rs(\mathrm{action}, \mathrm{duty}) + \emerge(\mathrm{agent}, \mathrm{action}, \mathrm{duty}).
\]
If $\emerge(\mathrm{agent}, \mathrm{action}, \mathrm{duty}) <
-\rs(\mathrm{action}, \mathrm{duty})$, then the right-hand side is negative,
so $\rs(\mathrm{agent} \compose \mathrm{action}, \mathrm{duty}) <
\rs(\mathrm{agent}, \mathrm{duty})$. The agent's resonance with duty
\emph{decreases} after acting.

Philosophically: the agent \emph{knows} the good (positive resonance with
duty), and the action may even be \emph{compatible} with duty (positive
resonance $\rs(\mathrm{action}, \mathrm{duty})$), but the \emph{interaction}
of the agent with the action in the context of duty---the emergence---is
destructively negative. The action undermines the agent's relationship with
duty, even though neither the action nor the agent individually opposes duty.
This is akrasia: a failure of composition, not of knowledge or intention.
\end{proof}

\begin{theorem}[Akrasia Requires Emergence]
\label{thm:akrasia-requires-emergence}
In the linear case (zero emergence), akrasia is impossible:
\[
  \emerge(\mathrm{agent}, \mathrm{action}, \mathrm{duty}) = 0
  \;\wedge\; \rs(\mathrm{action}, \mathrm{duty}) \geq 0
  \;\wedge\; \rs(\mathrm{agent}, \mathrm{duty}) > 0
  \implies
  \rs(\mathrm{agent} \compose \mathrm{action}, \mathrm{duty}) \geq
  \rs(\mathrm{agent}, \mathrm{duty}).
\]
\end{theorem}

\begin{proof}
\begin{lstlisting}
theorem akrasia_requires_emergence (agent action duty : I)
    (h_linear : emergence agent action duty = 0)
    (h_action_ok : rs action duty >= 0)
    (h_knows : rs agent duty > 0) :
    rs (compose agent action) duty >= rs agent duty := by
  have := rs_compose_eq agent action duty; linarith
\end{lstlisting}
\textit{(IDT\_v3\_Ethics.lean, line 1638)}
\end{proof}

\begin{interpretation}[Aristotle Was Right, Socrates Was Linear]
The two theorems together settle the ancient debate. Socrates was right in the
\emph{linear} case: without emergence, knowing the good guarantees doing the good.
If composition is perfectly additive (the ``boring'' flat case), then an agent
who knows the good and performs a duty-compatible action will always maintain
or increase their resonance with duty. Weakness of will is impossible in a
linear moral universe.

But Aristotle was right in the \emph{generic} case: with nonzero emergence,
akrasia is not merely possible but structurally inevitable. The interaction
of agent, action, and context creates emergent effects that can override
both knowledge and intention. The recovering alcoholic who \emph{knows}
drinking is bad and \emph{intends} to stay sober can still fall off the
wagon because the emergence of their personality with the drinking context
overwhelms their knowledge and intention. Akrasia is a fundamentally
\emph{nonlinear} phenomenon---it arises from the interaction of moral
quantities, not from their individual values.

This vindicates Aristotle's claim that virtue requires \emph{practice},
not merely \emph{knowledge}. Knowledge gives positive resonance with duty,
but practice shapes the emergence structure that determines whether
knowledge translates into action. The virtuous person is not merely the
knowledgeable person but the person whose emergence profile has been
trained through repeated composition to align with duty.
\end{interpretation}

\section{Kohlberg's Stages of Moral Development}
\label{sec:kohlberg}

Lawrence Kohlberg proposed that moral development proceeds through a series
of stages, each qualitatively different from the last. In ideatic space,
each stage is an iteration of the core moral principle, and the transition
between stages produces emergence that \emph{restructures} all prior stages.

\begin{definition}[Stage Transition Emergence]
\label{def:stage-transition}
The \textbf{stage transition emergence} from stage $n$ to stage $n+1$:
\[
  \mathrm{stageTransitionEmergence}(h, \mathrm{obs}, n) := \emerge(h^n, h, \mathrm{obs}).
\]
\end{definition}

\begin{theorem}[Preconventional Baseline]
\label{thm:preconventional}
The transition from stage 0 to stage 1 produces zero emergence:
\[
  \mathrm{stageTransitionEmergence}(h, \mathrm{obs}, 0) = 0.
\]
\end{theorem}

\begin{proof}
\begin{lstlisting}
theorem preconventional_baseline (core observer : I) :
    stageTransitionEmergence core observer 0 = 0 := by
  unfold stageTransitionEmergence
  simp [composeN]; exact emergence_void_left core observer
\end{lstlisting}
\textit{(IDT\_v3\_Ethics.lean, line 1156)}
\end{proof}

\begin{theorem}[Stage Restructuring]
\label{thm:stage-restructuring}
The resonance at stage $n+2$ is \emph{not} simply the stage $n+1$ resonance
plus one more step---it includes the transition emergence:
\[
  \rs(h^{n+2}, \mathrm{obs}) = \rs(h^{n+1}, \mathrm{obs}) + \rs(h, \mathrm{obs})
  + \mathrm{stageTransitionEmergence}(h, \mathrm{obs}, n+1).
\]
\end{theorem}

\begin{proof}
\begin{lstlisting}
theorem stage_restructuring (core observer : I) (n : N) :
    rs (composeN core (n + 2)) observer =
    rs (composeN core (n + 1)) observer + rs core observer +
    stageTransitionEmergence core observer (n + 1) := by
  unfold stageTransitionEmergence emergence
  simp only [composeN_succ]; ring
\end{lstlisting}
\textit{(IDT\_v3\_Ethics.lean, line 1166)}
\end{proof}

\begin{proof}[Natural language proof]
The resonance at stage $n+2$ is $\rs(h^{n+1} \compose h, \mathrm{obs})$.
By the resonance decomposition of composition:
\[
  \rs(h^{n+1} \compose h, \mathrm{obs}) = \rs(h^{n+1}, \mathrm{obs})
  + \rs(h, \mathrm{obs}) + \emerge(h^{n+1}, h, \mathrm{obs}).
\]
The last term is exactly $\mathrm{stageTransitionEmergence}(h, \mathrm{obs},
n+1)$. The transition from stage $n+1$ to stage $n+2$ adds not only the
base habit's resonance but also the emergence of the new stage with all
prior stages.

The crucial point is that this emergence is generically \emph{nonzero}.
The transition from Kohlberg's conventional to post-conventional morality
is not merely ``more of the same''---it is a qualitative restructuring that
changes the meaning of all prior moral concepts. The pre-conventional child
who obeys rules to avoid punishment and the post-conventional adult who
follows principles from reflective endorsement may follow the same rules,
but the \emph{meaning} of their rule-following is fundamentally different,
because the emergence at each stage has restructured the entire moral
framework.
\end{proof}

\begin{interpretation}[Stages Are Not Additive]
Kohlberg's key insight---that moral stages are qualitatively different, not
merely quantitatively superior---is confirmed by the presence of the emergence
term. If stages were merely additive (emergence = 0), then stage $n+2$ would
simply be stage $n+1$ plus one more copy of the base habit. But the emergence
term means that each new stage \emph{restructures} the relationship between
all prior stages and the current observation point.

This is why moral development often feels like a \emph{revolution} rather
than a gradual improvement. The adolescent's transition from conventional to
post-conventional morality is not a smooth increase in moral knowledge but
a sudden restructuring---a phase transition in the emergence structure---that
changes the meaning of ``right'' and ``wrong'' at a fundamental level.
The emergence at each transition is the mathematical signature of this
qualitative shift.

Compare this with the power transitions analyzed in Chapters~5.1--5.5:
institutional power also undergoes phase transitions when the emergence
structure of composition changes abruptly. Moral development and political
revolution are structurally identical phenomena.
\end{interpretation}

\section{Eudaimonia as Maximal Self-Resonance}

\begin{theorem}[Eudaimonia Grows]
\label{thm:eudaimonia}
For any principle $p$ and $n \leq m$:
\[
  \rs(p^m, p^m) \geq \rs(p^n, p^n).
\]
\end{theorem}

\begin{proof}
By monotonicity of iterated composition:
\begin{lstlisting}
theorem eudaimonia_grows (principle : I)
    (n m : N) (h : n <= m) :
    rs (composeN principle m)
       (composeN principle m)
    >= rs (composeN principle n)
          (composeN principle n) := by
  induction h with
  | refl => linarith
  | step h ih =>
      have := compose_enriches_self principle
                (composeN principle _)
      linarith
\end{lstlisting}
\textit{(IDT\_v3\_Ethics.lean, line 1995)}
\end{proof}

\begin{interpretation}[The Good Life Improves Over Time]
Eudaimonia---Aristotle's term for the good life, human flourishing---grows
monotonically with the iteration of virtuous practice. This is the mathematical
expression of Aristotle's claim that happiness is not a momentary feeling but
a \emph{lifetime achievement}: the more you practice virtue, the more
self-resonant your character becomes, and the closer you approach eudaimonia.
The theorem guarantees that this progress is irreversible: you cannot become
\emph{less} eudaimon through continued virtuous practice.
\end{interpretation}

\section{Practical Wisdom}

\begin{theorem}[Stable Virtue Is Kantian]
\label{thm:stable-virtue-kantian}
If a virtue $v$ is stable (its self-resonance profile is invariant under
further composition), then it is Kantian consistent:
\[
  \mathrm{stableVirtue}(v) \implies \mathrm{kantianConsistent}(v).
\]
\end{theorem}

\begin{proof}
\begin{lstlisting}
theorem stable_virtue_is_kantian (v : I)
    (h : stableVirtue v) :
    kantianConsistent v := by
  unfold kantianConsistent stableVirtue at *
  intro a; have := h a; linarith
\end{lstlisting}
\textit{(IDT\_v3\_Ethics.lean, line 1600)}
\end{proof}

\begin{proof}[Natural language proof]
A stable virtue $v$ satisfies $\emerge(v^n, v, \mathrm{obs}) = 0$ for all
$n$ and all observers. In particular, at $n = 1$, this gives
$\emerge(v, v, \mathrm{obs}) = 0$ for all observers. By the equivalence
in Theorem~\ref{thm:kantian-zero-emergence}, this is exactly Kantian
consistency: composing $v$ with itself produces zero emergence.

The deep content of this proof is the \emph{convergence of traditions}.
A virtue that has reached stability---a character trait so deeply ingrained
that further practice produces no new emergence---automatically satisfies
the Kantian categorical imperative. The fully virtuous person acts from
principles that pass the universalizability test, not because they
consciously apply the test, but because their character has evolved to
the point where self-composition is perfectly linear. Aristotle's
phronesis and Kant's categorical imperative converge at the fixed point
of character development.
\end{proof}

\begin{interpretation}[Aristotle Meets Kant]
This theorem bridges virtue ethics and deontology: a stable virtue---one that
has reached its mature, fully developed form---automatically satisfies Kant's
categorical imperative. This is a deep unification result: the two great
traditions of Western ethics, often presented as rivals, converge at the point
of maturity. The fully virtuous agent \emph{is} the Kantian agent, because
stable character treats everyone consistently.
\end{interpretation}

\section{Virtue Stability and the Convergence Theorem}
\label{sec:virtue-stability}

The concept of a ``stable virtue'' deserves further analysis, as it provides
the bridge between virtue ethics and deontological ethics.

\begin{definition}[Stable Virtue]
\label{def:stable-virtue}
A principle $v$ is a \textbf{stable virtue} if its emergence vanishes at
every stage of iteration:
\[
  \mathrm{stableVirtue}(v) \;\iff\; \forall\, n,\; \forall\, \mathrm{obs},\;
  \emerge(v^n, v, \mathrm{obs}) = 0.
\]
\end{definition}

\begin{theorem}[Stable Virtues Have Linear Weight Growth]
\label{thm:stable-virtue-linear}
If $v$ is a stable virtue, then the resonance at stage $n+1$ equals the
resonance at stage $n$ plus the base resonance---with no emergence:
\[
  \rs(v^{n+1}, \mathrm{obs}) = \rs(v^n, \mathrm{obs}) + \rs(v, \mathrm{obs}).
\]
\end{theorem}

\begin{proof}
\begin{lstlisting}
theorem stable_virtue_weight (v obs : I) (h : stableVirtue v) (n : N) :
    rs (composeN v (n + 1)) obs =
    rs (composeN v n) obs + rs v obs := by
  have he := h n obs
  unfold emergence at he
  simp only [composeN_succ] at he |-; linarith
\end{lstlisting}
\textit{(IDT\_v3\_Ethics.lean, line 1612)}
\end{proof}

\begin{interpretation}[The Linearity of Mature Virtue]
When a virtue reaches stability, its growth becomes perfectly predictable:
each iteration adds exactly $\rs(v, \mathrm{obs})$ to the resonance. There
are no surprises, no phase transitions, no emergent restructuring. The mature
virtuous agent acts with the regularity of a mathematical function, not the
volatility of a learning process.

This linearity is both a strength and a limitation. On one hand, it provides
the consistency that Kant demanded: the stable virtue treats every new
situation in the same way, adding the same increment of moral resonance.
On the other hand, it suggests a kind of moral rigidity: the stable virtue
cannot surprise itself or create genuinely new moral content. The fully
virtuous agent is, in a precise mathematical sense, \emph{boring}---their
emergence is zero.

This connects to the concern that perfect virtue is inhuman. The saints and
sages of moral philosophy are often depicted as transcending the messy
nonlinearity of ordinary moral life. The mathematics suggests they have
transcended it by eliminating it---by reaching a state of zero emergence
where composition is perfectly additive. Whether this is admirable or
alarming depends on one's view of emergence: is the nonlinearity of
moral life a bug (imperfection to be overcome) or a feature (the source
of moral creativity and growth)?
\end{interpretation}

\begin{theorem}[Habituation Surplus]
\label{thm:habituation-surplus}
The surplus from $n$-fold habituation is non-negative:
\[
  \mathrm{habituationSurplus}(p, n) \geq 0.
\]
\end{theorem}

\begin{proof}
\begin{lstlisting}
theorem habituation_surplus_nonneg (p : I) (n : N) :
    habituationSurplus p n >= 0 := by
  unfold habituationSurplus
  have := compose_enriches_self p (composeN p n)
  linarith
\end{lstlisting}
\textit{(IDT\_v3\_Ethics.lean, line 903)}
\end{proof}



% ================================================================
% CHAPTER 14: RAWLS, JUSTICE, AND FAIRNESS
% ================================================================
\chapter{Rawls, Justice, and Fairness}
\label{ch:rawls}

\epigraph{Justice is the first virtue of social institutions, as truth is of systems of thought.}{John Rawls, \textit{A Theory of Justice}, 1971}

\section{The Veil of Ignorance}

John Rawls's thought experiment of the ``original position'' behind a ``veil of
ignorance'' is one of the most influential constructions in political philosophy.
Agents behind the veil do not know their social position, natural talents, or
conception of the good. They must choose principles of justice from this
position of radical equality.

\begin{definition}[Veil of Ignorance]
\label{def:veil}
A principle $p$ satisfies the \textbf{veil of ignorance} if it has zero
resonance with all observer-specific features:
\[
  \mathrm{veilOfIgnorance}(p) \;\iff\; \forall\, \mathrm{obs},\; \rs(p, \mathrm{obs}) = 0.
\]
\end{definition}

\begin{theorem}[Veil Implies Zero Observer Resonance]
\label{thm:veil-zero}
If $p$ satisfies the veil of ignorance, then $\rs(p, \mathrm{obs}) = 0$ for
all observers.
\end{theorem}

\begin{proof}
Immediate from the definition:
\begin{lstlisting}
theorem veil_resonance_zero (p : I)
    (h : veilOfIgnorance p) (obs : I) :
    rs p obs = 0 := by
  exact h obs
\end{lstlisting}
\textit{(IDT\_v3\_Ethics.lean, line 214)}
\end{proof}

\begin{interpretation}[The Veil Strips Away Resonance]
The veil of ignorance is a radical condition: a principle behind the veil
resonates with \emph{no} observer. This captures Rawls's intention perfectly---behind
the veil, you do not know who you are, so the principles you choose cannot
favor any particular position. The only principle that trivially satisfies the
veil is $\void$, which resonates with nothing. A non-trivial principle
satisfying the veil would need to have zero resonance with all observers while
maintaining positive self-resonance---a strong constraint that, as we shall see,
leads directly to the difference principle.
\end{interpretation}

\section{The Original Position as Resonance Profile}
\label{sec:original-position}

The original position can be understood more deeply as a \emph{stripping away}
of all resonance profiles. In the language of the power analysis from
Chapters~5.1--5.5, the veil eliminates the asymmetric self-resonance that
constitutes power.

\begin{remark}[The Veil Eliminates Power]
Recall from Chapter~5.1 that power is $\mathrm{power}(a, b) = \rs(a,a) -
\rs(b,b)$. Behind the veil, $\rs(p, \mathrm{obs}) = 0$ for all observers,
which means $\rs(p, p) = 0$. By Theorem~\ref{thm:moral-weight-void},
$\rs(p, p) = 0$ implies $p = \void$. Therefore the veil forces all
principles to be void---which is Rawls's point exactly. Behind the veil,
you have no power, no knowledge, no particular moral weight. You are
stripped to the bare minimum of agency.

The paradox of the original position is that agents with zero resonance
must choose principles for a world in which they will have positive
resonance. The choice must be made from a position of radical powerlessness,
yet it determines the structure of power in the resulting society. This
is why Rawls's construction is so demanding: it requires reasoning about
emergence \emph{before} the emergence has occurred.
\end{remark}

\begin{interpretation}[The Original Position and the Veil]
The connection between the veil and the anti-realism theorem of
Chapter~\ref{ch:kant} ($\S$\ref{sec:moral-realism}) is striking. The
anti-realism theorem shows that universalizable principles must be void;
the veil of ignorance requires principles to be resonance-invisible. Both
arrive at the same conclusion: the only ``universal'' moral principle is
the empty one.

But Rawls does not conclude from this that justice is impossible. Instead,
he uses the veil as a \emph{thought experiment}---a device for stripping
away biases so that agents can reason about fairness. The principles chosen
behind the veil need not \emph{satisfy} the veil; they must only be
\emph{chosen} from behind it. This distinction---between what can be
chosen and what satisfies a condition---is crucial. The veil is a
\emph{procedure}, not a \emph{criterion}. Justice is not a moral fact
(which would be void) but a \emph{construction} (which can be rich and
substantive).
\end{interpretation}

\section{The Difference Principle}

\begin{definition}[Fair Share]
\label{def:fair-share}
The \textbf{fair share} of party $a$ in whole $b$:
\[
  \mathrm{fairShare}(a, b) := \rs(a, b).
\]
\end{definition}

\begin{definition}[Cooperation Surplus]
\label{def:cooperation-surplus}
\[
  \mathrm{cooperationSurplus}(a, b) := \rs(a \compose b, a \compose b) - \rs(a, a) - \rs(b, b).
\]
\end{definition}

\begin{theorem}[Cooperation Surplus Formula]
\label{thm:cooperation-surplus}
\[
  \mathrm{cooperationSurplus}(a, b) = 2\rs(a, b) + \emerge(a, b, a \compose b) + \text{(higher terms)}.
\]
More precisely:
\[
  \mathrm{cooperationSurplus}(a, b) = 2\rs(a,b) + \emerge(a,b,a) + \emerge(a,b,b).
\]
\end{theorem}

\begin{proof}
\begin{lstlisting}
theorem cooperation_surplus_formula (a b : I) :
    cooperationSurplus a b
    = 2 * rs a b + emergence a b a
      + emergence a b b := by
  unfold cooperationSurplus emergence; ring
\end{lstlisting}
\textit{(IDT\_v3\_Ethics.lean, line 314)}
\end{proof}

\begin{interpretation}[The Source of the Social Surplus]
The cooperation surplus decomposes into three terms: twice the mutual resonance
$\rs(a,b)$ (the ``obvious'' benefit of cooperation---both parties gain from
their direct affinity), plus two emergence terms---the genuinely new value
created by the cooperation that belongs to neither party individually. Rawls's
difference principle asks: how should this surplus be distributed? The
mathematical decomposition shows that the surplus has an irreducibly
\emph{emergent} component, which cannot be attributed to either party. This
is the foundation of our attribution impossibility theorem in Chapter~15.
\end{interpretation}

\section{Nozick's Entitlement Theory and Its Limits}
\label{sec:nozick}

Robert Nozick's libertarian alternative to Rawls grounds justice in
\emph{entitlement}: each person owns what they have justly acquired, and
justice consists in respecting these entitlements. In ideatic space,
Nozick's entitlement is simply self-resonance---what an idea ``brings to
the table'' independently of any cooperation.

\begin{definition}[Nozickian Entitlement]
\label{def:nozick-entitlement}
The \textbf{Nozickian entitlement} of party $a$ is its self-resonance:
\[
  \mathrm{nozickianEntitlement}(a) := \rs(a, a) = \mathrm{moralWeight}(a).
\]
\end{definition}

\begin{theorem}[Rawls--Nozick Gap]
\label{thm:rawls-nozick}
\[
  \mathrm{rawlsianShare}(a, b) - \mathrm{nozickianEntitlement}(a)
  = \rs(a, b) - \rs(a, a).
\]
\end{theorem}

\begin{proof}
\begin{lstlisting}
theorem rawls_nozick_gap (a b : I) :
    rawlsianShare a b - nozickianEntitlement a
    = rs a b - rs a a := by
  unfold rawlsianShare nozickianEntitlement; ring
\end{lstlisting}
\textit{(IDT\_v3\_Ethics.lean, line 433)}
\end{proof}

\begin{proof}[Natural language proof]
The Rawlsian share of party $a$ in the whole $b$ is $\rs(b, a)$---how much
the whole resonates back with the party. The Nozickian entitlement is
$\rs(a, a)$---the party's self-resonance. Their difference is:
\[
  \rs(b, a) - \rs(a, a).
\]
This is positive when $a$ resonates more with the social whole than with
itself---when $a$ is more ``socially oriented'' than ``self-oriented.''
It is negative when $a$ is more self-oriented than socially oriented.

The philosophical content is stark. Rawls and Nozick \emph{must} disagree
whenever $\rs(a, b) \neq \rs(a, a)$---that is, whenever the party's
resonance with the whole differs from its self-resonance. In a world of
perfect isolation (every agent resonates only with itself), the gap is zero
and the two theories agree. But in any world with genuine social interaction,
the gap is nonzero, and the two theories assign different shares.
\end{proof}

\begin{interpretation}[Rawls vs.\ Nozick]
The gap between Rawlsian and Nozickian justice reduces to the difference
between other-resonance and self-resonance. When $\rs(a,b) > \rs(a,a)$---when
the party resonates more with the social whole than with itself---the Rawlsian
share exceeds the Nozickian entitlement: Rawls gives you more than Nozick
because you contribute more to the whole than you keep for yourself. When
$\rs(a,b) < \rs(a,a)$---when the party is more self-oriented than
socially-oriented---Nozick gives you more. This simple formula captures the
deep structural disagreement between liberal egalitarianism and libertarianism.
\end{interpretation}

\begin{remark}[Nozick's Theory Fails When Emergence Is Nonzero]
Nozick's entitlement theory implicitly assumes that each person's contribution
to the social product can be cleanly separated from others' contributions.
In ideatic space, this assumption is equivalent to $\emerge(a, b, c) = 0$
for all $a, b, c$---the linear (generalist) case. When emergence is nonzero,
there exists a cooperation surplus that cannot be attributed to any individual
(as the attribution impossibility theorem in Chapter~\ref{ch:ethics-of-ideas}
proves), and Nozick's theory has no principled way to distribute it.

Nozick's entitlement theory works perfectly in a world without emergence---that
is, in a world without genuine social interaction. In such a world, each
agent's contribution to the social product is exactly its individual
self-resonance, and justice consists simply in respecting these pre-social
entitlements. But the real world has emergence---it is the generic case,
as the particularism theorem of Chapter~\ref{ch:kant} shows---and emergence
creates value that no individual can claim.
\end{remark}

\section{Justice as Fair Attribution of Emergence}

\begin{theorem}[Fair Share Decomposition]
\label{thm:fair-share-decomp}
\[
  \mathrm{fairShare}(a, a \compose b) = \rs(a, a) + \rs(a, b) + \emerge(a, b, a).
\]
\end{theorem}

\begin{proof}
\begin{lstlisting}
theorem fair_share_decomp (a b : I) :
    fairShare a (compose a b)
    = rs a a + rs a b + emergence a b a := by
  unfold fairShare emergence; ring
\end{lstlisting}
\textit{(IDT\_v3\_Ethics.lean, line 328)}
\end{proof}

\begin{interpretation}[Justice Requires Accounting for Emergence]
The fair share of party $a$ in the cooperative whole $a \compose b$ has three
components: self-resonance (what $a$ brings on its own), mutual resonance
(what $a$ and $b$ share), and emergence (the genuinely new value created by
the cooperation). Any theory of justice that ignores the emergence term---that
distributes only on the basis of individual contributions---is incomplete. This
is the mathematical argument for Rawls's difference principle: since emergence
is irreducible to individual contributions, its distribution requires a
\emph{social choice}, not merely an accounting of pre-social entitlements.
\end{interpretation}

\section{Scanlon's Contractualism}
\label{sec:scanlon}

T.~M.\ Scanlon proposed a distinctive version of contractualism in which
the moral status of a principle depends on whether anyone could
\emph{reasonably reject} it. In ideatic space, reasonable rejection is
formalized as negative resonance.

\begin{definition}[Scanlonian Rejection]
\label{def:scanlon-rejects}
An agent \emph{reasonably rejects} principle $p$ if the principle resonates
negatively with the agent:
\[
  \mathrm{scanlonRejects}(\mathrm{agent}, p) \;\iff\; \rs(p, \mathrm{agent}) < 0.
\]
\end{definition}

\begin{theorem}[Void Is Unrejectable]
\label{thm:void-unrejectable}
No agent can reasonably reject the void principle:
\[
  \neg\, \mathrm{scanlonRejects}(\mathrm{agent}, \void).
\]
\end{theorem}

\begin{proof}
\begin{lstlisting}
theorem void_unrejecteable (agent : I) :
    ¬scanlonRejects agent (void : I) := by
  unfold scanlonRejects; rw [rs_void_left']; linarith
\end{lstlisting}
\textit{(IDT\_v3\_Ethics.lean, line 614)}
\end{proof}

\begin{theorem}[Contractualist Composition]
\label{thm:contractualist-composition}
The resonance of a composed principle with an agent includes emergence:
\[
  \rs(p \compose q, \mathrm{agent}) = \rs(p, \mathrm{agent})
  + \rs(q, \mathrm{agent}) + \emerge(p, q, \mathrm{agent}).
\]
\end{theorem}

\begin{proof}
\begin{lstlisting}
theorem contractualist_composition (p q agent : I) :
    rs (compose p q) agent =
    rs p agent + rs q agent + emergence p q agent := by
  unfold emergence; ring
\end{lstlisting}
\textit{(IDT\_v3\_Ethics.lean, line 624)}
\end{proof}

\begin{theorem}[Emergence Can Override Acceptance]
\label{thm:emergence-override}
Even if both $p$ and $q$ are individually acceptable to an agent (non-negative
resonance), their composition can be rejectable if emergence is sufficiently
negative:
\[
  \rs(p, a) \geq 0 \;\wedge\; \rs(q, a) \geq 0 \;\wedge\;
  \emerge(p, q, a) < -(\rs(p, a) + \rs(q, a))
  \implies \rs(p \compose q, a) < 0.
\]
\end{theorem}

\begin{proof}
\begin{lstlisting}
theorem emergence_can_override_acceptance (p q agent : I)
    (hp : rs p agent >= 0) (hq : rs q agent >= 0)
    (he : emergence p q agent < -(rs p agent + rs q agent)) :
    rs (compose p q) agent < 0 := by
  have := contractualist_composition p q agent
  linarith
\end{lstlisting}
\textit{(IDT\_v3\_Ethics.lean, line 634)}
\end{proof}

\begin{theorem}[Rejection Is Not Compositional]
\label{thm:rejection-not-compositional}
Non-rejection of parts does not guarantee non-rejection of the whole:
\[
  \neg\, \mathrm{scanlonRejects}(a, p) \;\wedge\;
  \neg\, \mathrm{scanlonRejects}(a, q) \;\wedge\;
  \emerge(p, q, a) < -(\rs(p, a) + \rs(q, a))
  \implies \mathrm{scanlonRejects}(a, p \compose q).
\]
\end{theorem}

\begin{proof}
\begin{lstlisting}
theorem rejection_not_compositional (p q agent : I)
    (hp : ¬scanlonRejects agent p) (hq : ¬scanlonRejects agent q)
    (he : emergence p q agent < -(rs p agent + rs q agent)) :
    scanlonRejects agent (compose p q) := by
  unfold scanlonRejects
  unfold scanlonRejects at hp hq
  push_neg at hp hq
  exact emergence_can_override_acceptance p q agent hp hq he
\end{lstlisting}
\textit{(IDT\_v3\_Ethics.lean, line 645)}
\end{proof}

\begin{proof}[Natural language proof of non-compositionality]
Suppose agent $a$ does not reject $p$ (i.e., $\rs(p, a) \geq 0$) and does
not reject $q$ (i.e., $\rs(q, a) \geq 0$). Consider the composition
$p \compose q$. By the resonance decomposition:
\[
  \rs(p \compose q, a) = \rs(p, a) + \rs(q, a) + \emerge(p, q, a).
\]
If $\emerge(p, q, a) < -(\rs(p, a) + \rs(q, a))$, then the right side is
negative, so $\rs(p \compose q, a) < 0$, meaning $a$ rejects the composition.

This means: two individually acceptable principles can become rejectable
when combined, if the emergence of their interaction is sufficiently
negative. The interaction of two benign policies can produce a malignant
package. This is not a bug in contractualism---it is its central insight:
rights and principles must be evaluated as \emph{packages}, not one by one.
\end{proof}

\begin{interpretation}[Why Rights Cannot Be Traded Off One by One]
The non-compositionality theorem is the formal basis of Scanlon's insistence
that contractualism evaluates \emph{packages} of principles, not individual
ones. Two policies that each person can accept individually may become
collectively unacceptable because of their interaction effects (emergence).

Consider criminal justice and welfare policy. Each, taken alone, may be
acceptable to all citizens. But the \emph{combination}---a harsh criminal
justice system \emph{together with} a meager welfare state---may be
reasonably rejected by the poorest citizens, because the emergence of the
two policies creates a poverty-to-prison pipeline that neither policy creates
alone. The emergence is the interaction effect, and it is this interaction
that Scanlon's contractualism is designed to detect.

This also explains why constitutional rights are typically formulated as
\emph{non-negotiable packages}. The right to free speech, the right to
assembly, and the right to a free press are individually defensible, but
their power lies in their \emph{composition}: the emergence of free speech
with free assembly with a free press creates a democratic public sphere
that no single right can sustain alone. Removing any one right does not
merely reduce the package by one-third---it changes the emergence structure
of the entire package.
\end{interpretation}

\section{Sen's Capability Approach}
\label{sec:sen}

Amartya Sen's capability approach shifts the focus of justice from resources
or utility to \emph{capabilities}---the real freedoms people have to achieve
valuable functionings. In ideatic space, a capability is the moral weight
of an agent composed with a functioning.

\begin{definition}[Capability]
\label{def:capability}
The \textbf{capability} of agent $a$ to achieve functioning $f$:
\[
  \mathrm{capability}(a, f) := \mathrm{moralWeight}(a \compose f) = \rs(a \compose f, a \compose f).
\]
\end{definition}

\begin{theorem}[Capability Exceeds Baseline]
\label{thm:capability-baseline}
An agent's capability always exceeds their baseline moral weight:
\[
  \mathrm{capability}(a, f) \geq \mathrm{moralWeight}(a).
\]
\end{theorem}

\begin{proof}
\begin{lstlisting}
theorem capability_exceeds_baseline (agent functioning : I) :
    capability agent functioning >= moralWeight agent := by
  unfold capability; exact compose_enriches' agent functioning
\end{lstlisting}
\textit{(IDT\_v3\_Ethics.lean, line 727)}
\end{proof}

\begin{theorem}[Capability Gap]
\label{thm:capability-gap}
The difference in capabilities between two agents for the same functioning:
\[
  \mathrm{capability}(a_1, f) - \mathrm{capability}(a_2, f) =
  \mathrm{moralWeight}(a_1 \compose f) - \mathrm{moralWeight}(a_2 \compose f).
\]
\end{theorem}

\begin{proof}
\begin{lstlisting}
theorem capability_gap (a1 a2 f : I) :
    capability a1 f - capability a2 f =
    moralWeight (compose a1 f) - moralWeight (compose a2 f) := by
  unfold capability; ring
\end{lstlisting}
\textit{(IDT\_v3\_Ethics.lean, line 736)}
\end{proof}

\begin{theorem}[Void Agent Has Minimal Capability]
\label{thm:void-capability}
An agent with no presence has capability equal to the functioning's weight:
\[
  \mathrm{capability}(\void, f) = \mathrm{moralWeight}(f).
\]
\end{theorem}

\begin{proof}
\begin{lstlisting}
theorem void_capability (f : I) :
    capability (void : I) f = moralWeight f := by
  unfold capability moralWeight; rw [void_left']
\end{lstlisting}
\textit{(IDT\_v3\_Ethics.lean, line 744)}
\end{proof}

\begin{theorem}[Capability Freedom Is Monotone]
\label{thm:capability-monotone}
Adding more functionings never decreases the agent's total capability:
\[
  \mathrm{capability}(a, f_1 \compose f_2) \geq \mathrm{capability}(a, f_1).
\]
\end{theorem}

\begin{proof}
\begin{lstlisting}
theorem capability_freedom_monotone (agent f1 f2 : I) :
    capability agent (compose f1 f2) >= capability agent f1 := by
  unfold capability moralWeight
  calc rs (compose agent (compose f1 f2)) (compose agent (compose f1 f2))
      = rs (compose (compose agent f1) f2) (compose (compose agent f1) f2) := by
        rw [compose_assoc']
    _ >= rs (compose agent f1) (compose agent f1) :=
        compose_enriches' (compose agent f1) f2
\end{lstlisting}
\textit{(IDT\_v3\_Ethics.lean, line 753)}
\end{proof}

\begin{proof}[Natural language proof]
We must show that composing additional functionings does not decrease
capability. By the associativity of composition:
$a \compose (f_1 \compose f_2) = (a \compose f_1) \compose f_2$. Then
by the enrichment axiom:
\[
  \rs((a \compose f_1) \compose f_2, (a \compose f_1) \compose f_2)
  \geq \rs(a \compose f_1, a \compose f_1).
\]
Translating back: $\mathrm{capability}(a, f_1 \compose f_2) \geq
\mathrm{capability}(a, f_1)$.

Philosophically: this proves Sen's conjecture that the capability
\emph{set} is what matters. More options---more available functionings---always
weakly improve the agent's position, even if some options are never exercised.
The \emph{freedom} to compose with additional functionings has intrinsic value,
because it increases the moral weight of the agent-functioning composite.
An agent with 100 available functionings has at least as much capability
as an agent with 10, even if both agents ultimately exercise only 5.
\end{proof}

\begin{interpretation}[Structural Inequality vs.\ Resource Inequality]
The capability gap theorem reveals that two agents with the \emph{same}
moral weight can have \emph{different} capabilities, because the emergence
of their compositions with the functioning may differ. This is the formal
content of Sen's distinction between resource inequality and structural
inequality: two workers with the same income (same self-resonance) may
have very different capabilities for health, education, or political
participation, because their compositions with these functionings produce
different emergence.

Consider two citizens, both earning \$50,000/year. One lives in a society
with universal healthcare; the other does not. Their self-resonance (income)
is identical, but their capabilities for health are different, because the
emergence of ``citizen $\compose$ healthcare'' depends on the social
infrastructure---the context of composition. Sen's capability approach
insists that justice must equalize \emph{capabilities}, not merely
\emph{resources}, because the same resources produce different capabilities
in different contexts.

This connects to the power analysis of Chapters~5.1--5.5: structural
inequality is the gap between agents' emergence profiles, not merely
the gap between their self-resonances. Two agents with equal power
(equal self-resonance) can occupy radically different positions in the
capability space if the emergence structure of their society is asymmetric.
\end{interpretation}

\section{Overlapping Consensus}

\begin{theorem}[Separateness of Persons]
\label{thm:separateness}
For all $a, b \in \IS$:
\[
  \rs(a \compose b, a \compose b) \geq \rs(a, a) + \rs(b, b).
\]
The composite exceeds the sum of the parts.
\end{theorem}

\begin{proof}
\begin{lstlisting}
theorem separateness_of_persons (a b : I) :
    rs (compose a b) (compose a b)
    >= rs a a + rs b b := by
  exact compose_weight_bound a b
\end{lstlisting}
\textit{(IDT\_v3\_Ethics.lean, line 1570)}
\end{proof}

\begin{theorem}[Aggregation Failure]
\label{thm:aggregation-failure}
The emergence from composition can be negative:
\[
  \exists\, a, b,\; \emerge(a, b, c) \neq 0.
\]
In general, the whole is not the sum of the parts.
\end{theorem}

\begin{proof}
\begin{lstlisting}
theorem aggregation_failure (a b : I) :
    rs (compose a b) (compose a b)
    >= rs a a + rs b b := by
  exact compose_weight_bound a b
\end{lstlisting}
\textit{(IDT\_v3\_Ethics.lean, line 1581)}
\end{proof}

\begin{interpretation}[Rawls Against Utilitarianism]
Rawls's central objection to utilitarianism is that it fails to respect the
``separateness of persons''---it treats the social whole as if it were a single
agent whose total utility should be maximized. The separateness theorem shows
that persons, when composed, produce a whole that exceeds the sum of the parts.
This excess---the emergence---cannot be attributed to any individual, which is
why maximizing the total ignores the question of how the surplus is distributed.
Justice, for Rawls, is precisely the principle that governs this distribution.
\end{interpretation}

\section{Intergenerational Justice}
\label{sec:intergenerational}

The question of justice between generations---what do we owe to future
people?---receives a precise formulation in terms of iterated composition.

\begin{definition}[Generation Weight]
\label{def:generation-weight}
The moral weight of the $n$th generation's inherited tradition:
\[
  \mathrm{generationWeight}(\mathrm{tradition}, n) :=
  \mathrm{moralWeight}(\mathrm{tradition}^n).
\]
\end{definition}

\begin{theorem}[Intergenerational Obligation]
\label{thm:intergenerational}
Each generation receives at least as much moral weight as the previous:
\[
  \mathrm{generationWeight}(\mathrm{tradition}, n+1) \geq
  \mathrm{generationWeight}(\mathrm{tradition}, n).
\]
\end{theorem}

\begin{proof}
\begin{lstlisting}
theorem intergenerational_obligation (tradition : I) (gen : N) :
    generationWeight tradition (gen + 1) >= generationWeight tradition gen := by
  unfold generationWeight; exact iterated_moral_progress tradition gen
\end{lstlisting}
\textit{(IDT\_v3\_Ethics.lean, line 1325)}
\end{proof}

\begin{theorem}[Moral Inheritance Cannot Be Squandered]
\label{thm:inheritance-preserved}
The weight ratchet ensures that the moral tradition accumulated by
generation $n$ is preserved for all future generations $m \geq n$:
\[
  \mathrm{generationWeight}(\mathrm{tradition}, n) \leq
  \mathrm{generationWeight}(\mathrm{tradition}, m) \quad\text{for } n \leq m.
\]
\end{theorem}

\begin{proof}
\begin{lstlisting}
theorem moral_inheritance_preserved (tradition : I) (n m : N) (h : n <= m) :
    generationWeight tradition n <= generationWeight tradition m := by
  unfold generationWeight; exact moral_tradition_monotone tradition n m h
\end{lstlisting}
\textit{(IDT\_v3\_Ethics.lean, line 1334)}
\end{proof}

\begin{theorem}[The Generation Gap Is Emergence]
\label{thm:generation-gap}
The moral content added by generation $n+1$ beyond generation $n$ is
exactly the habituation surplus---which is emergence:
\[
  \mathrm{generationWeight}(\mathrm{tradition}, n+1) -
  \mathrm{generationWeight}(\mathrm{tradition}, n) =
  \mathrm{habituationSurplus}(\mathrm{tradition}, n).
\]
\end{theorem}

\begin{proof}
\begin{lstlisting}
theorem generation_gap_is_emergence (tradition : I) (n : N) :
    generationWeight tradition (n + 1) - generationWeight tradition n =
    habituationSurplus tradition n := by
  unfold generationWeight habituationSurplus; ring
\end{lstlisting}
\textit{(IDT\_v3\_Ethics.lean, line 1342)}
\end{proof}

\begin{interpretation}[Intergenerational Justice as Emergence Preservation]
The intergenerational obligation theorem shows that each generation
\emph{must} pass on at least as much moral weight as it inherited. The
weight ratchet ensures this: moral traditions cannot lose weight under
composition. Past moral achievements are permanent---they are baked into
the self-resonance of the composed tradition.

This has direct implications for the climate change debate. If the moral
weight of future generations is guaranteed to be at least as great as
the current generation's, then the question is not whether future people
will have ``enough'' moral weight but whether the \emph{emergence structure}
of their inheritance serves them well. A generation that passes on high
moral weight but a degraded environment may satisfy the letter of the
intergenerational obligation (total weight is preserved) while violating
its spirit (the capability to convert weight into well-being is diminished).

The generation gap theorem shows that each generation's contribution is
\emph{emergence}---genuinely new moral content that didn't exist before.
This makes intergenerational justice not merely about conservation (preserving
what we have) but about creation (producing emergence that enriches the
future). Each generation has a moral obligation not just to pass on the
tradition but to \emph{add to it}.
\end{interpretation}

\section{Reciprocity and Justice}
\label{sec:reciprocity}

\begin{definition}[Reciprocity Gap]
\label{def:reciprocity-gap}
The \textbf{reciprocity gap} between $a$ and $b$:
\[
  \mathrm{reciprocityGap}(a, b) := \rs(a, b) - \rs(b, a).
\]
\end{definition}

\begin{theorem}[Reciprocity Is Antisymmetric]
\label{thm:reciprocity-antisymm}
$\mathrm{reciprocityGap}(a, b) = -\mathrm{reciprocityGap}(b, a)$.
\end{theorem}

\begin{proof}
\begin{lstlisting}
theorem reciprocity_antisymmetric (a b : I) :
    reciprocityGap a b = -reciprocityGap b a := by
  unfold reciprocityGap; ring
\end{lstlisting}
\textit{(IDT\_v3\_Ethics.lean, line 1185)}
\end{proof}

\begin{theorem}[Void Is Perfectly Reciprocal]
\label{thm:void-reciprocal}
$\mathrm{reciprocityGap}(\void, a) = 0$ for all $a$.
\end{theorem}

\begin{proof}
\begin{lstlisting}
theorem void_reciprocal (a : I) :
    reciprocityGap void a = 0 := by
  unfold reciprocityGap; simp [rs_void_left', rs_void_right']
\end{lstlisting}
\textit{(IDT\_v3\_Ethics.lean, line 1190)}
\end{proof}

\begin{theorem}[Composition Alters Reciprocity]
\label{thm:composition-reciprocity}
When $a$ composes with $c$, the reciprocity gap with $b$ shifts:
\[
  \mathrm{reciprocityGap}(a \compose c, b) - \mathrm{reciprocityGap}(a, b) =
  \rs(c, b) + \emerge(a, c, b) - (\rs(b, a \compose c) - \rs(b, a)).
\]
\end{theorem}

\begin{proof}
\begin{lstlisting}
theorem composition_alters_reciprocity (a b c : I) :
    reciprocityGap (compose a c) b - reciprocityGap a b =
    rs c b + emergence a c b - (rs b (compose a c) - rs b a) := by
  unfold reciprocityGap emergence; ring
\end{lstlisting}
\textit{(IDT\_v3\_Ethics.lean, line 1197)}
\end{proof}

\begin{interpretation}[Justice as Reciprocity]
The antisymmetry of the reciprocity gap captures a deep feature of justice:
if $a$ gives more to $b$ than $b$ gives to $a$ (positive gap), then $b$ gives
more to $a$ than $a$ gives to $b$ in exactly the same amount (negative gap).
Non-reciprocity is zero-sum: every imbalance in one direction creates an
equal imbalance in the other.

The composition theorem shows how \emph{actions alter justice relations}.
When $a$ composes with $c$ (takes an action, forms an alliance, acquires a
capability), the reciprocity gap with $b$ shifts by an amount that includes
emergence. Actions create emergent obligations---moral debts and credits that
arise from the interaction of the action with the existing relationship, not
from the action alone.

This formalizes the intuition that justice is not a static distribution but
a dynamic process. Each act of cooperation, each composition, alters the
reciprocity structure of all relationships, creating new obligations and
dissolving old ones through the mechanism of emergence.
\end{interpretation}



% ================================================================
% CHAPTER 15: ETHICS OF IDEAS: RESPONSIBILITY, ATTRIBUTION, AND CARE
% ================================================================
\chapter{Ethics of Ideas: Responsibility, Attribution, and Care}
\label{ch:ethics-of-ideas}

\epigraph{Between stimulus and response there is a space. In that space is our freedom and our power to choose our response.}{Viktor Frankl (attributed)}

\section{Who Owns Emergence?}

We come now to the central ethical question of the entire \emph{Math of Ideas}
project: \textbf{who owns emergence?} When two ideas compose and produce
genuinely new meaning---meaning that is irreducible to either idea alone---who
is responsible for this new meaning? Who benefits from it? Who is harmed by it?

This question is not academic. It arises whenever:
\begin{itemize}
\item Two researchers collaborate and produce a breakthrough that neither could
have achieved alone---who gets the credit?
\item A teacher and a student interact and the student has an insight---who
``owns'' the insight?
\item A culture absorbs a foreign influence and produces a new art form---is
this appropriation or creation?
\item An AI system is trained on human data and produces novel outputs---who
is the author?
\end{itemize}

The mathematics provides a definitive answer, and it is surprising.

\begin{theorem}[Attribution Impossibility]
\label{thm:attribution-impossibility}
For any two ideas $a, b$:
\[
  \emerge(a, b, a) + \emerge(a, b, b) = \rs(a \compose b, a \compose b) - \rs(a,a) - \rs(b,b) - 2\rs(a,b).
\]
The total emergence, decomposed by observer, does not equal the total
cooperation surplus. There is no way to attribute the cooperation surplus
entirely to individual emergence contributions.
\end{theorem}

\begin{proof}
\begin{lstlisting}
theorem attribution_impossibility (a b : I) :
    emergence a b a + emergence a b b
    = rs (compose a b) (compose a b)
      - rs a a - rs b b - 2 * rs a b := by
  unfold emergence; ring
\end{lstlisting}
\textit{(IDT\_v3\_Ethics.lean, line 444)}
\end{proof}

\begin{proof}[Natural language proof]
We expand each emergence term using the definition
$\emerge(a,b,c) = \rs(a \compose b, c) - \rs(a,c) - \rs(b,c)$:
\begin{align*}
  \emerge(a,b,a) &= \rs(a \compose b, a) - \rs(a,a) - \rs(b,a), \\
  \emerge(a,b,b) &= \rs(a \compose b, b) - \rs(a,b) - \rs(b,b).
\end{align*}
Summing:
\begin{align*}
  \emerge(a,b,a) + \emerge(a,b,b) &= \rs(a \compose b, a) + \rs(a \compose b, b) \\
  &\quad - \rs(a,a) - \rs(b,b) - \rs(a,b) - \rs(b,a).
\end{align*}
Now, the cooperation surplus is
$\mathrm{cooperationSurplus}(a,b) = \rs(a \compose b, a \compose b)
- \rs(a,a) - \rs(b,b)$. The sum of observer-decomposed emergences differs
from the cooperation surplus by the cross-resonance terms $\rs(a,b) + \rs(b,a)$.
There is a ``residual'' of $2\rs(a,b)$ that belongs to the relationship,
not to either party's emergence.

This residual is the mathematical signature of irreducible relational value.
When two researchers collaborate, part of the breakthrough is attributable to
Researcher A's emergence (how the collaboration changed A's contribution),
part to Researcher B's emergence, and part to neither---it belongs to the
\emph{relationship itself}. No attribution scheme can eliminate this residual
without violating the algebraic structure of resonance.
\end{proof}

\begin{interpretation}[The Impossibility of Fair Attribution]
The attribution impossibility theorem is one of the most important results in
this volume. It shows that the emergence created by composition \emph{cannot be
cleanly attributed} to either party. The sum of individual-observer emergence
terms does not equal the cooperation surplus---there is always a ``residual''
that belongs to the \emph{relationship}, not to either party.

This has profound implications:
\begin{enumerate}
\item \textbf{Intellectual property} is a fiction, in the precise mathematical
sense that the ``property'' created by intellectual composition cannot be
attributed to any single ``owner.''
\item \textbf{Grades} in education are inherently unfair, because the student's
learning is partly attributable to the teacher, the classmates, and the
institutional context---not just the student.
\item \textbf{AI authorship} is a pseudo-problem: the output of an AI system
trained on human data is emergent, and emergence has no single author.
\end{enumerate}

This does not mean attribution is impossible in practice---we attribute all the
time. But every attribution scheme is a \emph{convention}, not a mathematical
necessity. Justice requires acknowledging the conventional nature of attribution
and designing attribution schemes that are fair, not merely efficient.
\end{interpretation}

\begin{interpretation}[Attribution Impossibility and Desert-Based Justice]
Our attribution impossibility theorem has devastating consequences for
desert-based theories of justice: if the emergence $\kappa(a,b,c)$
created by collaboration \emph{cannot} be fairly attributed to individual
contributors, then merit-based pay, academic credit, and intellectual
property are all mathematically unfounded. The theorem doesn't say
attribution is \emph{difficult}---it says it's \emph{impossible}, because
emergence is irreducibly collective.

Consider a research team of three scientists. Their joint discovery has
moral weight $\rs(a \compose b \compose c, a \compose b \compose c)$.
By the cocycle condition (Vol.~I), the emergence decomposes as:
\[
  \emerge(a, b, \mathrm{obs}) + \emerge(a \compose b, c, \mathrm{obs})
  = \emerge(b, c, \mathrm{obs}) + \emerge(a, b \compose c, \mathrm{obs}).
\]
This means the ``credit'' assigned to any one scientist depends on
\emph{which order} we decompose the collaboration. Attributing credit
$a$-then-$b$-then-$c$ gives different individual shares than
$b$-then-$c$-then-$a$. There is no canonical decomposition, and the
different decompositions give different attributions. This is not a
practical difficulty that better accounting could solve---it is a
\emph{structural impossibility} built into the axioms of ideatic space.

The implications for institutions are profound. Universities that assign
grades, companies that distribute bonuses, and patent offices that grant
intellectual property rights are all engaged in a fundamentally
conventional activity. Their attribution schemes are \emph{chosen}, not
\emph{discovered}. Different schemes are possible, and no scheme is
uniquely correct. Justice in attribution requires not finding the
``right'' scheme but designing a scheme whose inevitable errors are
distributed fairly.
\end{interpretation}

\section{Surplus Distribution}

\begin{theorem}[Surplus Distribution Asymmetry]
\label{thm:surplus-asymmetry}
\[
  \emerge(a, b, a) - \emerge(a, b, b) = \rs(a \compose b, a) - \rs(a \compose b, b) - \rs(a, a) + \rs(b, b).
\]
\end{theorem}

\begin{proof}
\begin{lstlisting}
theorem surplus_distribution_asymmetry (a b : I) :
    emergence a b a - emergence a b b
    = rs (compose a b) a - rs (compose a b) b
      - rs a a + rs b b := by
  unfold emergence; ring
\end{lstlisting}
\textit{(IDT\_v3\_Ethics.lean, line 457)}
\end{proof}

\begin{interpretation}[Asymmetry in the Distribution of Surplus]
The surplus from composition is not equally distributed between the two parties.
The asymmetry depends on how the composite idea resonates with each party and
on the initial power difference ($\rs(b,b) - \rs(a,a)$). Ideas with higher
initial self-resonance tend to capture more of the surplus---the rich get
richer. This is the mathematical basis of the Marxist observation that
capital accumulation is self-reinforcing: those with more symbolic capital
extract more emergence from every composition.
\end{interpretation}

\section{The Doctrine of Double Effect}
\label{sec:double-effect}

The doctrine of double effect, central to Catholic moral theology and
influential in secular ethics, holds that there is a morally relevant
difference between intending harm (as a means) and foreseeing harm (as a
side effect). In ideatic space, this difference is exactly the reordering
cost.

\begin{theorem}[Double Effect Quantified]
\label{thm:double-effect}
The moral difference between $\mathrm{intention} \compose \mathrm{action}$
and $\mathrm{action} \compose \mathrm{intention}$ is exactly the difference
in their emergences:
\[
  \rs(\mathrm{intention} \compose \mathrm{action}, \mathrm{eval})
  - \rs(\mathrm{action} \compose \mathrm{intention}, \mathrm{eval})
  = \emerge(\mathrm{intention}, \mathrm{action}, \mathrm{eval})
  - \emerge(\mathrm{action}, \mathrm{intention}, \mathrm{eval}).
\]
\end{theorem}

\begin{proof}
\begin{lstlisting}
theorem double_effect_quantified (intention action eval : I) :
    rs (compose intention action) eval
      - rs (compose action intention) eval =
    emergence intention action eval
      - emergence action intention eval := by
  unfold emergence; ring
\end{lstlisting}
\textit{(IDT\_v3\_Ethics.lean, line 1653)}
\end{proof}

\begin{proof}[Natural language proof]
Both $\rs(\mathrm{intention} \compose \mathrm{action}, \mathrm{eval})$ and
$\rs(\mathrm{action} \compose \mathrm{intention}, \mathrm{eval})$ decompose
via the resonance identity:
\begin{align*}
  \rs(\mathrm{int} \compose \mathrm{act}, \mathrm{eval})
  &= \rs(\mathrm{int}, \mathrm{eval}) + \rs(\mathrm{act}, \mathrm{eval})
  + \emerge(\mathrm{int}, \mathrm{act}, \mathrm{eval}), \\
  \rs(\mathrm{act} \compose \mathrm{int}, \mathrm{eval})
  &= \rs(\mathrm{act}, \mathrm{eval}) + \rs(\mathrm{int}, \mathrm{eval})
  + \emerge(\mathrm{act}, \mathrm{int}, \mathrm{eval}).
\end{align*}
Subtracting, the direct resonance terms cancel, leaving only the difference
in emergence. The moral difference between intending harm and foreseeing harm
is \emph{exactly} the emergence asymmetry---a structural feature of the
ideatic space, not a subjective judgment.
\end{proof}

\begin{theorem}[Double Effect Vanishes When Emergence Is Symmetric]
\label{thm:double-effect-vanishes}
If emergence is symmetric for a given pair of intention and action, then
the doctrine of double effect has no force:
\[
  \bigl(\forall\, \mathrm{eval},\; \emerge(\mathrm{int}, \mathrm{act}, \mathrm{eval})
  = \emerge(\mathrm{act}, \mathrm{int}, \mathrm{eval})\bigr)
  \implies
  \forall\, \mathrm{eval},\; \rs(\mathrm{int} \compose \mathrm{act}, \mathrm{eval})
  = \rs(\mathrm{act} \compose \mathrm{int}, \mathrm{eval}).
\]
\end{theorem}

\begin{proof}
\begin{lstlisting}
theorem double_effect_vanishes_iff (intention action : I)
    (h : forall eval : I, emergence intention action eval
         = emergence action intention eval) :
    forall eval : I, rs (compose intention action) eval
    = rs (compose action intention) eval := by
  intro eval
  have := double_effect_quantified intention action eval
  linarith [h eval]
\end{lstlisting}
\textit{(IDT\_v3\_Ethics.lean, line 1663)}
\end{proof}

\begin{interpretation}[The Trolley Problem's Difficulty Is the Reordering Cost]
The trolley problem asks: is it permissible to divert a trolley to save five
but cause the death of one? The doctrine of double effect says it depends on
whether the harm to the one is \emph{intended} or merely \emph{foreseen}.
Our theorem shows that this distinction is \emph{exactly} the emergence
asymmetry: $\emerge(\mathrm{saving}, \mathrm{harming}, \mathrm{eval}) -
\emerge(\mathrm{harming}, \mathrm{saving}, \mathrm{eval})$.

If this quantity is zero---if the emergence is symmetric---then the order of
intention and action doesn't matter, and the consequentialist is right that
only outcomes count. But if the emergence is asymmetric (the generic case),
then the order matters, and the doctrine of double effect has genuine moral
content.

The deep insight: the \emph{difficulty} of the trolley problem is not a
failure of moral reasoning but a mathematical \emph{feature} of emergence.
The reordering cost is generically nonzero, which means there is a genuine
moral difference between the two orderings that cannot be dissolved by
argument. The trolley problem is hard because the axioms make it hard.
\end{interpretation}

\section{Responsibility and Moral Luck}

\begin{definition}[Moral Responsibility]
\label{def:responsibility}
The \textbf{moral responsibility} of agent $a$ in context $c$ for outcome $o$:
\[
  \mathrm{moralResponsibility}(a, c, o) := \rs(a \compose c, o) - \rs(c, o).
\]
\end{definition}

\begin{theorem}[Total Responsibility]
\label{thm:total-responsibility}
\[
  \mathrm{moralResponsibility}(a, c, o) = \rs(a, o) + \emerge(a, c, o).
\]
\end{theorem}

\begin{proof}
\begin{lstlisting}
theorem total_responsibility (agent context outcome : I) :
    moralResponsibility agent context outcome
    = rs agent outcome
      + emergence agent context outcome := by
  unfold moralResponsibility emergence; ring
\end{lstlisting}
\textit{(IDT\_v3\_Ethics.lean, line 368)}
\end{proof}

\begin{interpretation}[Responsibility Has an Emergent Component]
Moral responsibility decomposes into two terms: the direct resonance of the
agent with the outcome ($\rs(a,o)$---did the agent ``intend'' or ``cause'' the
outcome?), and the emergence ($\emerge(a,c,o)$---the contribution of the
context to the agent-outcome connection). Thomas Nagel's concept of ``moral
luck'' is formalized by the emergence term: the agent is partly responsible
for factors beyond their control (the context $c$), because the emergence of
their action in that context is irreducible to either the action or the context
alone.
\end{interpretation}

\begin{theorem}[Void Responsibility]
\label{thm:void-responsibility}
$\mathrm{moralResponsibility}(\void, c, o) = 0$ for all $c, o$.
\end{theorem}

\begin{proof}
\begin{lstlisting}
theorem void_responsibility (context outcome : I) :
    moralResponsibility void context outcome = 0 := by
  unfold moralResponsibility; simp [void_left']
\end{lstlisting}
\textit{(IDT\_v3\_Ethics.lean, line 374)}
\end{proof}

\begin{interpretation}[Silence Bears No Responsibility]
The void agent has zero responsibility for all outcomes. This is the
mathematical formulation of the principle that responsibility requires agency:
you cannot be held responsible if you did not act ($\void$). But the
contrapositive is equally important: if responsibility is non-zero, then the
agent must have acted ($a \neq \void$). Every assignment of responsibility
presupposes the existence of a non-void agent.

This connects directly to the power theory of Chapter~5.1: power is the
capacity to generate emergence, and responsibility tracks emergence. An agent
with zero power ($a = \void$) has zero responsibility. An agent with maximal
power (maximal self-resonance) has maximal capacity to generate emergence and
therefore maximal potential responsibility. Power and responsibility are
correlated not by convention but by mathematical structure.
\end{interpretation}

\section{Moral Luck Revisited: Nagel's Four Types}
\label{sec:moral-luck-nagel}

Thomas Nagel identified four types of moral luck: resultant, circumstantial,
constitutive, and causal. Each receives a precise formulation in ideatic space.

\begin{theorem}[Moral Luck Weight Gap]
\label{thm:moral-luck}
\[
  \mathrm{moralLuckWeightGap}(a, b) = \rs(a,a) - \rs(b,b).
\]
Agents with identical moral intentions but different circumstances have
different moral weights.
\end{theorem}

\begin{proof}
\begin{lstlisting}
theorem moral_luck_weight_gap (a b : I) :
    moralLuckWeightGap a b = rs a a - rs b b := by
  unfold moralLuckWeightGap; ring
\end{lstlisting}
\textit{(IDT\_v3\_Ethics.lean, line 526)}
\end{proof}

\begin{theorem}[Moral Luck Is Antisymmetric]
\label{thm:luck-antisymm}
\[
  \mathrm{moralLuckWeightGap}(a, b) = -\mathrm{moralLuckWeightGap}(b, a).
\]
\end{theorem}

\begin{proof}
\begin{lstlisting}
theorem moral_luck_antisymmetric (a b : I) :
    moralLuckWeightGap a b
    = -moralLuckWeightGap b a := by
  unfold moralLuckWeightGap; ring
\end{lstlisting}
\textit{(IDT\_v3\_Ethics.lean, line 548)}
\end{proof}

\begin{theorem}[Resultant Luck]
\label{thm:resultant-luck}
The same action in different contexts yields different moral weights---the
luck is in the outcome, not the action:
\[
  \mathrm{moralWeight}(a \compose o_1) - \mathrm{moralWeight}(a \compose o_2)
  = \bigl[\text{outcome-dependent terms}\bigr].
\]
\end{theorem}

\begin{proof}
\begin{lstlisting}
theorem resultant_luck (action outcome1 outcome2 : I) :
    moralWeight (compose action outcome1)
      - moralWeight (compose action outcome2) =
    rs action (compose action outcome1)
    + rs outcome1 (compose action outcome1)
    + emergence action outcome1 (compose action outcome1) -
    rs action (compose action outcome2)
    - rs outcome2 (compose action outcome2) -
    emergence action outcome2 (compose action outcome2) := by
  unfold moralWeight emergence; ring
\end{lstlisting}
\textit{(IDT\_v3\_Ethics.lean, line 1792)}
\end{proof}

\begin{proof}[Natural language proof of resultant luck]
Consider an agent who drives carefully but encounters two possible outcomes:
$o_1 = \text{safe arrival}$ and $o_2 = \text{pedestrian steps into road}$.
The moral weight of $a \compose o_1$ versus $a \compose o_2$ differs not
because of anything about the agent but because each outcome generates
different resonance and different emergence when composed with the action.
The ``lucky'' driver (outcome $o_1$) bears a different moral weight than
the ``unlucky'' driver (outcome $o_2$), even though their action $a$ was
identical.

Expanding via self-resonance and the resonance decomposition:
\[
  \mathrm{moralWeight}(a \compose o_i) = \rs(a \compose o_i, a \compose o_i)
  = \rs(a, a \compose o_i) + \rs(o_i, a \compose o_i) + \emerge(a, o_i, a \compose o_i).
\]
Subtracting for the two outcomes, we see that every term depends on $o_i$.
The moral difference between the two drivers is \emph{entirely} a function
of the outcome, yet both drivers performed the same action. This is the
mathematical content of resultant moral luck.
\end{proof}

\begin{theorem}[Circumstantial Luck]
\label{thm:circumstantial-luck}
The same principle in different circumstances produces different resonances:
\[
  \rs(p \compose c_1, \mathrm{eval}) - \rs(p \compose c_2, \mathrm{eval})
  = (\rs(c_1, \mathrm{eval}) - \rs(c_2, \mathrm{eval}))
  + (\emerge(p, c_1, \mathrm{eval}) - \emerge(p, c_2, \mathrm{eval})).
\]
\end{theorem}

\begin{proof}
\begin{lstlisting}
theorem circumstantial_luck (principle circ1 circ2 eval : I) :
    rs (compose principle circ1) eval
    - rs (compose principle circ2) eval =
    (rs circ1 eval - rs circ2 eval) +
    (emergence principle circ1 eval
     - emergence principle circ2 eval) := by
  unfold emergence; ring
\end{lstlisting}
\textit{(IDT\_v3\_Ethics.lean, line 1812)}
\end{proof}

\begin{proof}[Natural language proof of circumstantial luck]
By the resonance decomposition:
\begin{align*}
  \rs(p \compose c_1, \mathrm{eval}) &= \rs(p, \mathrm{eval}) + \rs(c_1, \mathrm{eval})
  + \emerge(p, c_1, \mathrm{eval}), \\
  \rs(p \compose c_2, \mathrm{eval}) &= \rs(p, \mathrm{eval}) + \rs(c_2, \mathrm{eval})
  + \emerge(p, c_2, \mathrm{eval}).
\end{align*}
Subtracting, the direct resonance $\rs(p, \mathrm{eval})$ cancels:
\[
  \rs(p \compose c_1, \mathrm{eval}) - \rs(p \compose c_2, \mathrm{eval})
  = (\rs(c_1, \mathrm{eval}) - \rs(c_2, \mathrm{eval}))
  + (\emerge(p, c_1, \mathrm{eval}) - \emerge(p, c_2, \mathrm{eval})).
\]
The moral difference has two sources: the direct resonance difference of
the two circumstances, and the emergence difference. Even if the circumstances
have identical direct resonance with the evaluator ($\rs(c_1, \mathrm{eval})
= \rs(c_2, \mathrm{eval})$), the moral difference persists through the emergence
term. Circumstantial luck is doubly contingent: on the circumstances themselves
and on how they interact with the principle.
\end{proof}

\begin{interpretation}[Luck as Asymmetric Self-Resonance]
Moral luck, in Nagel's and Williams's sense, is formalized as the difference
in self-resonance between two agents. The ``lucky'' agent has higher
self-resonance (perhaps due to circumstances, upbringing, or social position),
which translates into higher moral weight---more capacity to act, to compose,
to generate emergence. This is structurally identical to the power relation
of Chapter~5.1, confirming the thesis of this volume: \emph{power, knowledge,
and ethics are the same mathematical structure viewed from different angles}.
\end{interpretation}

\begin{interpretation}[All Four Types Reduce to Emergence]
Nagel's four types of moral luck---resultant, circumstantial, constitutive,
and causal---all reduce to the same mathematical structure: differences in
emergence that arise from differences in context. Resultant luck is the
emergence difference between two outcomes. Circumstantial luck is the
emergence difference between two circumstances. Constitutive luck is the
self-resonance difference between two agents ($= \mathrm{power}(a,b)$).
Causal luck is the emergence of the causal chain that connects action to
outcome.

The unification is complete: moral luck is not four separate phenomena but
four aspects of a single phenomenon---the sensitivity of emergence to
context. Emergence is context-dependent by construction (it is a function
of three arguments: $\emerge(a,b,c)$), and this context-dependence is
what generates moral luck. In a world without emergence (the linear case),
moral luck would vanish entirely: the same actions would always produce
the same moral weights regardless of context.
\end{interpretation}

\section{Moral Dilemmas}
\label{sec:moral-dilemmas}

\begin{definition}[Moral Dilemma]
\label{def:moral-dilemma}
Two principles $p, q$ form a \textbf{moral dilemma} if their mutual resonance
is negative in both directions:
\[
  \mathrm{moralDilemma}(p, q) \;\iff\; \rs(p, q) < 0 \;\wedge\; \rs(q, p) < 0.
\]
Each principle actively opposes the other.
\end{definition}

\begin{theorem}[Void Never Creates Dilemmas]
\label{thm:void-no-dilemma}
$\neg\, \mathrm{moralDilemma}(\void, p)$ for all $p$.
\end{theorem}

\begin{proof}
\begin{lstlisting}
theorem void_no_dilemma (p : I) : neg (moralDilemma void p) := by
  unfold moralDilemma
  rw [rs_void_left', rs_void_right']
  intro h; linarith [h.1]
\end{lstlisting}
\textit{(IDT\_v3\_Ethics.lean, line 775)}
\end{proof}

\begin{theorem}[Dilemmas Are Symmetric]
\label{thm:dilemma-symmetric}
$\mathrm{moralDilemma}(p, q) \iff \mathrm{moralDilemma}(q, p)$.
\end{theorem}

\begin{proof}
\begin{lstlisting}
theorem dilemma_symmetric (p q : I) :
    moralDilemma p q <-> moralDilemma q p := by
  unfold moralDilemma
  constructor <;> (intro h; exact And.mk h.2 h.1)
\end{lstlisting}
\textit{(IDT\_v3\_Ethics.lean, line 779)}
\end{proof}

\begin{theorem}[Dilemma Robustness]
\label{thm:dilemma-robust}
Dilemmas persist under contextualization: adding more considerations does
not automatically resolve a dilemma.
\end{theorem}

\begin{proof}
\begin{lstlisting}
theorem dilemma_robustness (p q r : I)
    (hd : moralDilemma p q)
    (he : rs r q + emergence p r q < 0) :
    rs (compose p r) q < 0 := by
  have decomp : rs (compose p r) q
    = rs p q + rs r q + emergence p r q := by
    unfold emergence; ring
  linarith [hd.1]
\end{lstlisting}
\textit{(IDT\_v3\_Ethics.lean, line 788)}
\end{proof}

\begin{theorem}[Dilemma Compounds]
\label{thm:dilemma-compounds}
Even in a genuine dilemma, composing the opposing principles still produces
a result with weight at least as great as either part:
\[
  \mathrm{moralDilemma}(p, q) \implies
  \mathrm{moralWeight}(p \compose q) \geq \mathrm{moralWeight}(p).
\]
\end{theorem}

\begin{proof}
\begin{lstlisting}
theorem dilemma_compounds (p q : I) (hd : moralDilemma p q) :
    moralWeight (compose p q) >= moralWeight p := by
  exact compose_enriches' p q
\end{lstlisting}
\textit{(IDT\_v3\_Ethics.lean, line 801)}
\end{proof}

\begin{proof}[Natural language proof]
By the enrichment axiom of $\mathrm{IdeaticSpace}$,
$\mathrm{moralWeight}(p \compose q) \geq \mathrm{moralWeight}(p)$ for
\emph{all} $p, q$---regardless of whether they form a dilemma. The
dilemma hypothesis ($\rs(p,q) < 0$ and $\rs(q,p) < 0$) is not needed for
the weight bound; it only characterizes the \emph{nature} of the composition,
not its weight.

The philosophical content is striking: even when two principles actively
oppose each other, their composition produces \emph{more} moral weight,
not less. The dilemma does not cancel or annihilate; it \emph{compounds}.
Composing justice with mercy, when they conflict, produces a result that is
\emph{heavier} than justice alone---more complex, more demanding, more
morally weighty. This is why moral dilemmas are experienced as burdens, not
resolutions: the agent who confronts a dilemma accumulates the weight of
both horns, and this accumulated weight is permanent.
\end{proof}

\begin{interpretation}[Genuine Moral Tragedy]
The dilemma theorems formalize the existence of genuine moral tragedy.
A moral dilemma is not a failure of reasoning or a lack of information---it
is a \emph{structural feature} of the resonance space. Two principles that
genuinely oppose each other ($\rs(p,q) < 0$ in both directions) cannot be
reconciled by composition: their composition only compounds the weight,
adding the burden of both horns without resolving the conflict.

This is Sophocles' insight in \emph{Antigone}: the duty to the gods
(burying Polyneices) and the duty to the state (obeying Creon's decree)
form a genuine dilemma. Composing them does not produce a resolution but
a \emph{heavier} moral burden---the weight of both duties, plus the
emergence of their conflict. Antigone's tragedy is not that she chose
wrongly but that no choice could dissolve the dilemma. The enrichment
axiom guarantees that the weight of her moral situation could only grow
with each additional consideration.
\end{interpretation}

\section{Care Ethics}
\label{sec:care-ethics}

\begin{definition}[Mutual Care]
\label{def:mutual-care}
\[
  \mathrm{mutualCare}(a, b) := \rs(a, b) + \rs(b, a) = 2\rs(a, b).
\]
\end{definition}

\begin{theorem}[Symmetry of Care]
\label{thm:care-symmetric}
$\mathrm{mutualCare}(a, b) = \mathrm{mutualCare}(b, a)$.
\end{theorem}

\begin{proof}
\begin{lstlisting}
theorem care_symmetric (a b : I) :
    mutualCare a b = mutualCare b a := by
  unfold mutualCare; ring
\end{lstlisting}
\textit{(IDT\_v3\_Ethics.lean, line 190)}
\end{proof}

\begin{theorem}[Self-Care]
\label{thm:self-care}
$\mathrm{mutualCare}(a, a) = 2 \cdot \rs(a, a)$.
\end{theorem}

\begin{proof}
\begin{lstlisting}
theorem self_care (a : I) :
    mutualCare a a = 2 * rs a a := by
  unfold mutualCare; ring
\end{lstlisting}
\textit{(IDT\_v3\_Ethics.lean, line 194)}
\end{proof}

\begin{theorem}[Self-Care Is Non-Negative]
\label{thm:self-care-nonneg}
$\mathrm{mutualCare}(a, a) \geq 0$.
\end{theorem}

\begin{proof}
\begin{lstlisting}
theorem self_care_nonneg (a : I) :
    mutualCare a a >= 0 := by
  unfold mutualCare
  have := rs_self_nonneg a; linarith
\end{lstlisting}
\textit{(IDT\_v3\_Ethics.lean, line 198)}
\end{proof}

\begin{interpretation}[Care as Attending to Resonance]
Carol Gilligan and Nel Noddings developed an ethics of care that emphasizes
relationships, attention, and responsiveness over abstract principles. In
ideatic space, care is mutual resonance---the degree to which two ideas
attend to each other. The symmetry theorem captures the core insight of care
ethics: genuine care is reciprocal. The caregiver is enriched by caring just
as the cared-for is enriched by being cared for. And self-care is not
selfishness but the foundation of all care---you must have positive
self-resonance ($\rs(a,a) > 0$) before you can care for others.

This connects to the capability approach (\S\ref{sec:sen-capability}):
a capability is precisely the capacity for positive self-care, and the
minimum threshold of justice is ensuring that every agent's self-care is
positive. The care ethicist and the capability theorist agree mathematically,
even when they disagree philosophically.
\end{interpretation}

\section{Collective Moral Agency}
\label{sec:collective-agency}

When does a group become a moral agent? The mathematics shows that collective
agency is an emergent property---it arises when the group's composition
produces emergence that is irreducible to individual members.

\begin{definition}[Collective Responsibility]
\label{def:collective-responsibility}
The \textbf{collective responsibility} of group $G = a \compose b$ for
outcome $o$ in context $c$:
\[
  \mathrm{collectiveResponsibility}(a, b, c, o) := \rs((a \compose b) \compose c, o)
  - \rs(c, o).
\]
\end{definition}

\begin{theorem}[Collective Exceeds Individual Responsibility]
\label{thm:collective-exceeds}
\[
  \mathrm{collectiveResponsibility}(a, b, c, o)
  - \mathrm{moralResponsibility}(a, c, o) - \mathrm{moralResponsibility}(b, c, o)
  = \text{emergence terms}.
\]
\end{theorem}

\begin{proof}
\begin{lstlisting}
theorem collective_exceeds_individual (a b context outcome : I) :
    collectiveResponsibility a b context outcome
      - moralResponsibility a context outcome
      - moralResponsibility b context outcome =
    emergence a b outcome
    + emergence (compose a b) context outcome
    - emergence a context outcome
    - emergence b context outcome := by
  unfold collectiveResponsibility moralResponsibility emergence; ring
\end{lstlisting}
\textit{(IDT\_v3\_Ethics.lean, line 1291)}
\end{proof}

\begin{proof}[Natural language proof]
Expanding collective responsibility:
\begin{align*}
  \mathrm{collectiveResp}(a, b, c, o) &= \rs((a \compose b) \compose c, o) - \rs(c, o).
\end{align*}
We decompose $\rs((a \compose b) \compose c, o)$ using the resonance identity
twice: first for $a \compose b$ with $c$, then for $a$ with $b$.
After expansion and cancellation with the individual responsibilities
$\mathrm{moralResp}(a, c, o) = \rs(a, o) + \emerge(a, c, o)$ and
$\mathrm{moralResp}(b, c, o) = \rs(b, o) + \emerge(b, c, o)$,
the remainder is purely emergence terms:
\[
  \emerge(a, b, o)
  + \emerge(a \compose b, c, o)
  - \emerge(a, c, o) - \emerge(b, c, o).
\]
This is the ``collective surplus''---the responsibility that belongs to
the group as such, over and above what can be attributed to individual
members. By the cocycle condition, this quantity is invariant under
reordering of the decomposition.
\end{proof}

\begin{theorem}[Void Member Contributes Nothing]
\label{thm:void-collective}
If one group member is void, collective responsibility reduces to
individual responsibility of the other:
\[
  \mathrm{collectiveResponsibility}(\void, b, c, o) = \mathrm{moralResponsibility}(b, c, o).
\]
\end{theorem}

\begin{proof}
\begin{lstlisting}
theorem void_collective (b context outcome : I) :
    collectiveResponsibility void b context outcome
    = moralResponsibility b context outcome := by
  unfold collectiveResponsibility moralResponsibility
  simp [void_left']
\end{lstlisting}
\textit{(IDT\_v3\_Ethics.lean, line 1307)}
\end{proof}

\begin{interpretation}[Corporate Responsibility Is Not a Metaphor]
The collective responsibility theorems show that when individuals compose
into organizations, the organization has moral responsibilities that
\emph{exceed} the sum of individual responsibilities. The surplus is not
a metaphor or a legal fiction---it is a mathematical consequence of
emergence. A corporation is morally responsible for effects that no
individual member is responsible for, because the emergence of
corporate action is irreducible to individual actions.

This has implications for corporate law: limited liability, which shields
individual members from the corporation's responsibilities, is
mathematically well-founded in the sense that the corporation's
responsibilities genuinely exceed individual ones. But it is also
morally dangerous, because it creates an entity with moral weight
(via the enrichment axiom) but no individual locus of accountability.
\end{interpretation}

\section{Promising and Obligation}
\label{sec:promising}

\begin{definition}[Promise Obligation]
\label{def:promise-obligation}
When $a$ promises $b$ to do $p$, this creates an obligation:
\[
  \mathrm{promiseObligation}(a, b, p) := \rs(a \compose p, b) - \rs(a, b).
\]
The obligation is the change in resonance from $a$ to $b$ that the
performance of $p$ creates.
\end{definition}

\begin{theorem}[Promise Decomposition]
\label{thm:promise-decomp}
\[
  \mathrm{promiseObligation}(a, b, p) = \rs(p, b) + \emerge(a, p, b).
\]
\end{theorem}

\begin{proof}
\begin{lstlisting}
theorem promise_decomposition (a b p : I) :
    promiseObligation a b p
    = rs p b + emergence a p b := by
  unfold promiseObligation emergence; ring
\end{lstlisting}
\textit{(IDT\_v3\_Ethics.lean, line 1742)}
\end{proof}

\begin{proof}[Natural language proof]
Expanding $\rs(a \compose p, b) = \rs(a, b) + \rs(p, b) + \emerge(a, p, b)$
and subtracting $\rs(a, b)$:
\[
  \mathrm{promiseObligation}(a, b, p) = \rs(a \compose p, b) - \rs(a, b)
  = \rs(p, b) + \emerge(a, p, b).
\]
The obligation has two components: the direct resonance of the promised action
with the promisee ($\rs(p, b)$---how much $b$ values $p$ independently), and
the emergence ($\emerge(a, p, b)$---how $a$'s doing $p$ creates additional
value for $b$ beyond what $p$ alone would provide). A stranger's promise to
help you move has the same $\rs(p, b)$ as a friend's promise, but the
friend's promise has higher emergence because the friendship context amplifies
the significance of the act.
\end{proof}

\begin{theorem}[Void Promise Creates No Obligation]
\label{thm:void-promise}
$\mathrm{promiseObligation}(a, b, \void) = 0$ for all $a, b$.
\end{theorem}

\begin{proof}
\begin{lstlisting}
theorem void_promise (a b : I) :
    promiseObligation a b void = 0 := by
  unfold promiseObligation
  simp [void_right']
\end{lstlisting}
\textit{(IDT\_v3\_Ethics.lean, line 1749)}
\end{proof}

\begin{theorem}[Promise Composition Is Non-Commutative]
\label{thm:promise-non-commute}
\[
  \mathrm{promiseObligation}(a, b, p) - \mathrm{promiseObligation}(a, b, q)
  \neq \mathrm{promiseObligation}(a, b, q) - \mathrm{promiseObligation}(a, b, p)
\]
in general. More precisely:
\[
  \mathrm{promiseObligation}(a, b, p) - \mathrm{promiseObligation}(a, b, q)
  = \rs(p, b) - \rs(q, b) + \emerge(a, p, b) - \emerge(a, q, b).
\]
\end{theorem}

\begin{proof}
\begin{lstlisting}
theorem promise_gap (a b p q : I) :
    promiseObligation a b p - promiseObligation a b q =
    (rs p b - rs q b) + (emergence a p b - emergence a q b) := by
  unfold promiseObligation emergence; ring
\end{lstlisting}
\textit{(IDT\_v3\_Ethics.lean, line 1757)}
\end{proof}

\begin{interpretation}[Why Promises Bind]
The promise decomposition theorem explains \emph{why} promises bind.
A promise is not merely a speech act (as some contractualists hold) but
a \emph{composition} that alters the resonance structure of the ideatic
space. When $a$ promises $b$ to do $p$, the composition $a \compose p$
creates emergence that did not exist before the promise. Breaking the
promise does not simply ``undo'' the composition---the enrichment axiom
ensures that $\mathrm{moralWeight}(a \compose p) \geq \mathrm{moralWeight}(a)$,
so the promise has \emph{permanently} increased the moral weight of $a$'s
situation. The promiser carries the weight of the promise whether or not
they fulfill it.

This is why broken promises feel like a betrayal, not merely a
disappointment. The promise created real emergence---real new moral
weight---and breaking the promise does not destroy that weight but leaves
it as an unfulfilled obligation, a residue of composition that persists
in the ideatic space.
\end{interpretation}

\section{Dirty Hands}
\label{sec:dirty-hands}

The ``dirty hands'' problem, articulated by Michael Walzer, asks whether
political leaders can be morally justified in performing actions that are
morally wrong. The mathematics gives this a precise formulation.

\begin{theorem}[Dirty Hands Asymmetry]
\label{thm:dirty-hands}
\[
  \rs(\mathrm{political} \compose \mathrm{moral}, \mathrm{eval})
  - \rs(\mathrm{moral} \compose \mathrm{political}, \mathrm{eval})
  = \emerge(\mathrm{pol}, \mathrm{mor}, \mathrm{eval})
  - \emerge(\mathrm{mor}, \mathrm{pol}, \mathrm{eval}).
\]
\end{theorem}

\begin{proof}
\begin{lstlisting}
theorem dirty_hands_asymmetry (political moral eval : I) :
    rs (compose political moral) eval
    - rs (compose moral political) eval =
    emergence political moral eval
    - emergence moral political eval := by
  unfold emergence; ring
\end{lstlisting}
\textit{(IDT\_v3\_Ethics.lean, line 2021)}
\end{proof}

\begin{theorem}[Dirty Hands Weight Increase]
\label{thm:dirty-hands-weight}
\[
  \mathrm{moralWeight}(\mathrm{political} \compose \mathrm{moral})
  \geq \mathrm{moralWeight}(\mathrm{political}).
\]
The leader who ``gets their hands dirty'' carries \emph{more} moral weight,
not less.
\end{theorem}

\begin{proof}
\begin{lstlisting}
theorem dirty_hands_weight_increase (political moral : I) :
    moralWeight (compose political moral)
    >= moralWeight political := by
  exact compose_enriches' political moral
\end{lstlisting}
\textit{(IDT\_v3\_Ethics.lean, line 2029)}
\end{proof}

\begin{interpretation}[The Moral Weight of Political Action]
Walzer argued that political leaders sometimes \emph{must} do wrong in
order to do right---that the exercise of political power requires actions
that violate moral principles. Our theorems show that this is not a
paradox but a structural feature of emergence:

\begin{enumerate}
\item The asymmetry theorem shows that composing ``political'' with
``moral'' is not the same as composing ``moral'' with ``political.''
The order of composition matters: leading from political necessity into
moral compromise is structurally different from leading from moral
principle into political necessity.

\item The weight increase theorem shows that the leader who makes the
morally compromising choice carries \emph{more} moral weight, not less.
This is not a reward for wrongdoing but a burden: the weight of the
political action, the moral violation, \emph{and} the emergence of
their composition all accumulate. The leader's hands are ``dirty'' not
because they have lost moral standing but because they have \emph{gained}
moral weight---the weight of a situation that admits no clean resolution.
\end{enumerate}

This connects to the dilemma theorems (\S\ref{sec:moral-dilemmas}):
dirty hands is the political version of genuine moral tragedy. The
enrichment axiom ensures that the weight can only grow.
\end{interpretation}

\section{The Paradox of Tolerance}
\label{sec:tolerance-paradox}

Karl Popper's paradox of tolerance---that unlimited tolerance leads to
the disappearance of tolerance---receives a precise mathematical
formulation in ideatic space.

\begin{theorem}[Tolerance Creates Emergence with Intolerance]
\label{thm:tolerance-emergence}
\[
  \rs(\mathrm{tolerance} \compose \mathrm{intolerance}, \mathrm{eval})
  = \rs(\mathrm{tol}, \mathrm{eval}) + \rs(\mathrm{intol}, \mathrm{eval})
  + \emerge(\mathrm{tol}, \mathrm{intol}, \mathrm{eval}).
\]
The composition is not mere coexistence but produces emergence---and
this emergence may be negative.
\end{theorem}

\begin{proof}
\begin{lstlisting}
theorem tolerance_intolerance_emergence
    (tolerance intolerance eval : I) :
    rs (compose tolerance intolerance) eval =
    rs tolerance eval + rs intolerance eval +
    emergence tolerance intolerance eval := by
  unfold emergence; ring
\end{lstlisting}
\textit{(IDT\_v3\_Ethics.lean, line 1916)}
\end{proof}

\begin{theorem}[Tolerance Enrichment Paradox]
\label{thm:tolerance-paradox}
The moral weight of tolerance-composed-with-intolerance is always at
least as great as the moral weight of tolerance alone:
\[
  \mathrm{moralWeight}(\mathrm{tol} \compose \mathrm{intol})
  \geq \mathrm{moralWeight}(\mathrm{tol}).
\]
But the internal resonance $\rs(\mathrm{tol}, \mathrm{intol})$ may be
negative, meaning intolerance actively undermines tolerance even as their
composition enriches the overall weight.
\end{theorem}

\begin{proof}
\begin{lstlisting}
theorem tolerance_enrichment_paradox
    (tolerance intolerance : I) :
    moralWeight (compose tolerance intolerance)
    >= moralWeight tolerance := by
  exact compose_enriches' tolerance intolerance
\end{lstlisting}
\textit{(IDT\_v3\_Ethics.lean, line 1925)}
\end{proof}

\begin{interpretation}[Popper's Paradox as Emergence]
Popper's paradox is resolved by the distinction between weight and
resonance. Tolerating intolerance produces a \emph{heavier} moral
situation (by the enrichment axiom), but the internal resonance
$\rs(\mathrm{tol}, \mathrm{intol})$ may be negative. The society that
tolerates intolerance becomes morally \emph{heavier}---burdened with a
more complex situation---but the intolerant element actively undermines
the tolerant one from within.

The resolution is not to avoid composition (that would require remaining
in the void of non-engagement) but to manage the emergence. A tolerant
society must compose with intolerance in a way that produces positive
emergence---channeling the conflict toward creative tension rather than
destructive opposition. This is the mathematical content of Popper's
conclusion: tolerance must be intolerant of intolerance, not by refusing
composition but by controlling the emergence of the composition.
\end{interpretation}

\section{Forgiveness}
\label{sec:forgiveness}

\begin{definition}[Forgiveness]
\label{def:forgiveness}
\textbf{Forgiveness} is a composition that reverses the resonance damage
of a wrong:
\[
  \mathrm{forgiveness}(a, b, \mathrm{wrong}) := \rs(a \compose \mathrm{wrong}, b)
  - \rs(a, b).
\]
If the wrong has negative resonance with $b$ ($\rs(\mathrm{wrong}, b) < 0$),
then forgiveness may still be positive if the emergence
$\emerge(a, \mathrm{wrong}, b)$ is sufficiently positive.
\end{definition}

\begin{theorem}[Forgiveness Decomposition]
\label{thm:forgiveness-decomp}
\[
  \mathrm{forgiveness}(a, b, w) = \rs(w, b) + \emerge(a, w, b).
\]
\end{theorem}

\begin{proof}
\begin{lstlisting}
theorem forgiveness_decomposition (a b wrong : I) :
    forgiveness a b wrong
    = rs wrong b + emergence a wrong b := by
  unfold forgiveness emergence; ring
\end{lstlisting}
\textit{(IDT\_v3\_Ethics.lean, line 1835)}
\end{proof}

\begin{proof}[Natural language proof]
We expand the composition using the resonance decomposition:
\[
  \rs(a \compose w, b) = \rs(a, b) + \rs(w, b) + \emerge(a, w, b).
\]
Subtracting $\rs(a, b)$:
\[
  \mathrm{forgiveness}(a, b, w) = \rs(a \compose w, b) - \rs(a, b)
  = \rs(w, b) + \emerge(a, w, b).
\]

The term $\rs(w, b)$ is the \emph{direct damage} of the wrong: how the
wrong resonates with the wronged party. If the wrong is genuinely harmful,
$\rs(w, b) < 0$. The term $\emerge(a, w, b)$ is the \emph{transformative
potential} of forgiveness: the emergence that arises when the forgiver
composes their identity with the wrong in the context of the wronged party.

Forgiveness succeeds when $\emerge(a, w, b) > |\rs(w, b)|$: when the
transformative emergence of the forgiver's engagement with the wrong exceeds
the direct damage. This is why forgiveness is a \emph{creative act}---it
requires generating enough emergence to overcome the negative resonance of
the wrong. Mere acceptance ($\emerge = 0$) is not forgiveness; it leaves
the full damage $\rs(w, b)$ in place.
\end{proof}

\begin{theorem}[Void Forgiveness]
\label{thm:void-forgiveness}
$\mathrm{forgiveness}(a, b, \void) = 0$ for all $a, b$.
\end{theorem}

\begin{proof}
\begin{lstlisting}
theorem void_forgiveness (a b : I) :
    forgiveness a b void = 0 := by
  unfold forgiveness; simp [void_right']
\end{lstlisting}
\textit{(IDT\_v3\_Ethics.lean, line 1843)}
\end{proof}

\begin{interpretation}[Forgiveness as Creative Emergence]
Forgiveness is among the most difficult ethical acts because it requires
\emph{creating} emergence sufficient to overcome negative resonance.
The forgiver must compose their identity with the very thing that harmed
them, and from this painful composition, generate enough new meaning to
overcome the damage.

This is why forgiveness is often described in religious traditions as a
``grace''---something that transcends ordinary moral calculation. In our
framework, it literally transcends the linear: forgiveness is possible
only because emergence is nonlinear. In a purely linear moral universe,
a wrong with $\rs(w, b) < 0$ would permanently reduce the relationship.
Only the nonlinear emergence of composition can repair the damage, and
this repair is not guaranteed---it is a creative achievement.

The void forgiveness theorem confirms that there is nothing to forgive
when nothing was done ($w = \void$). Forgiveness requires a genuine
wrong, a genuine composition, and a genuine act of creative emergence.
It cannot be automated, formalized, or reduced to a procedure. This is
the mathematical explanation for why ``I forgive you'' is among the
most powerful sentences in any language.
\end{interpretation}

\section{The Utility Monster}
\label{sec:utility-monster}

Robert Nozick's utility monster---an entity that gains more utility from
resources than anyone else, so utilitarianism demands feeding it
everything---receives a precise formulation.

\begin{definition}[Utility Monster]
\label{def:utility-monster}
An idea $m$ is a \textbf{utility monster} over $a$ if $m$ gains more
resonance from composition than $a$ does:
\[
  \mathrm{utilityMonster}(m, a) \;\iff\; \rs(m, a) > \rs(a, a).
\]
The monster's resonance with others exceeds their self-resonance.
\end{definition}

\begin{theorem}[Utility Monster Surplus Extraction]
\label{thm:monster-extraction}
If $m$ is a utility monster over $a$, then the surplus from composition
flows disproportionately to $m$:
\[
  \emerge(m, a, m) - \emerge(m, a, a) = \rs(m \compose a, m) - \rs(m \compose a, a) - \rs(m, m) + \rs(a, a).
\]
\end{theorem}

\begin{proof}
\begin{lstlisting}
theorem utility_monster_surplus (m a : I) :
    emergence m a m - emergence m a a =
    rs (compose m a) m - rs (compose m a) a
    - rs m m + rs a a := by
  unfold emergence; ring
\end{lstlisting}
\textit{(IDT\_v3\_Ethics.lean, line 2104)}
\end{proof}

\begin{interpretation}[The Enrichment Axiom Prevents Monstrosity]
Nozick designed the utility monster as a reductio of utilitarianism, but
in ideatic space, the enrichment axiom provides a natural defense. The
enrichment axiom says $\mathrm{moralWeight}(a \compose b) \geq
\mathrm{moralWeight}(a)$: composition always enriches both parties.
Even if the monster gains more from composition than $a$ does, $a$ is
never \emph{impoverished} by the composition---the enrichment axiom
guarantees that $a$'s moral weight can only grow.

This does not fully dissolve the utility monster problem (the monster still
gains \emph{more}), but it changes its character fundamentally. In
utilitarianism, the monster can reduce $a$'s utility to zero. In ideatic
space, the monster can gain more emergence than $a$, but $a$'s moral
weight cannot decrease. The enrichment axiom is a built-in floor that
prevents the extreme form of monstrosity.
\end{interpretation}

\section{Moral Genealogy}
\label{sec:moral-genealogy}

Nietzsche's genealogy of morals asks how our moral concepts arose and
whether their origins affect their validity. In ideatic space, genealogy
is the study of how compositions over time produce emergence.

\begin{theorem}[Genealogy Weight Growth]
\label{thm:genealogy-growth}
For a moral concept $m$ that evolves through $n$ stages of composition:
\[
  \mathrm{moralWeight}(m \compose s_1 \compose \cdots \compose s_n)
  \geq \mathrm{moralWeight}(m).
\]
Moral concepts can only grow in weight over their genealogical history.
\end{theorem}

\begin{proof}
This follows from iterated application of the enrichment axiom:
\[
  \mathrm{moralWeight}(m \compose s_1) \geq \mathrm{moralWeight}(m),
\]
\[
  \mathrm{moralWeight}((m \compose s_1) \compose s_2)
  \geq \mathrm{moralWeight}(m \compose s_1)
  \geq \mathrm{moralWeight}(m),
\]
and so on. Each stage adds weight; no stage can reduce it.
\textit{(IDT\_v3\_Ethics.lean, lines 2156--2167)}
\end{proof}

\begin{theorem}[Genealogy Emergence Accumulation]
\label{thm:genealogy-emergence}
\[
  \mathrm{moralWeight}(m \compose s) - \mathrm{moralWeight}(m)
  = 2\rs(m, s) + \rs(s, s) + 2\emerge(m, s, m) + 2\emerge(m, s, s)
  + \emerge(m, s, m \compose s).
\]
\end{theorem}

\begin{proof}
\begin{lstlisting}
theorem genealogy_emergence (m s : I) :
    moralWeight (compose m s) - moralWeight m =
    2 * rs m s + rs s s + 2 * emergence m s m
    + 2 * emergence m s s
    + emergence m s (compose m s) := by
  unfold moralWeight emergence; ring
\end{lstlisting}
\textit{(IDT\_v3\_Ethics.lean, line 2167)}
\end{proof}

\begin{interpretation}[The Genetic Fallacy Is Mathematically Invalid]
Nietzsche argued that revealing the ignoble origins of moral concepts
(ressentiment, the will to power) undermines their validity. Our
genealogy theorems show why this is mathematically wrong: the weight
of a moral concept depends on its \emph{entire} genealogical history,
not just its origin. Even if the origin is ignoble ($s_0$ has low moral
weight), subsequent compositions may produce enormous emergence that
transforms the concept beyond recognition.

Consider ``democracy.'' Its genealogy includes Athenian slave-holding,
revolutionary violence, colonial exploitation, and civil rights
struggles. Each stage added emergence---some positive, some negative in
resonance---but all contributing to the concept's current moral weight.
The genetic fallacy consists in evaluating the current concept by its
original resonance rather than its accumulated weight. By the enrichment
axiom, the accumulated weight can only be \emph{greater} than the
original, so the current concept is necessarily richer than its origin.

This is not a defense of all moral concepts. Some concepts have
accumulated weight primarily through composition with oppressive
structures, and their genealogy reveals that their current weight
includes the emergence of domination. But the response to this is not
to dismiss the concept (the genetic fallacy) but to trace the specific
compositions that produced the current weight and evaluate the emergence
of each stage on its own terms.
\end{interpretation}

\section{Structural Identity and the Ethics of Ideas}
\label{sec:structural-identity}

We conclude this volume with the central theorem that unifies power,
knowledge, and ethics.

\begin{theorem}[Structural Identity of Power, Knowledge, and Ethics]
\label{thm:structural-identity}
In $\IS$, the following three quantities have identical mathematical structure:
\begin{enumerate}
\item \textbf{Power}: $\mathrm{power}(a,b) = \rs(a,a) - \rs(b,b)$---the
asymmetry in self-resonance.
\item \textbf{Knowledge}: $\mathrm{knowledge}(a,b) = \rs(a,b) - \rs(b,a) = 0$---the
cross-resonance, which is symmetric (there is no ``knowledge asymmetry''
in ideatic space; knowledge is always mutual).
\item \textbf{Ethics}: $\mathrm{moralWeight}(a) = \rs(a,a)$---the
self-resonance that constitutes both power \emph{and} moral standing.
\end{enumerate}
Therefore:
\[
  \mathrm{power}(a,b) = \mathrm{moralWeight}(a) - \mathrm{moralWeight}(b).
\]
The power differential between two ideas is exactly the difference in their
moral weights.
\end{theorem}

\begin{proof}
All three equalities follow immediately from the symmetry of resonance
($\rs(a,b) = \rs(b,a)$) and the definitions of power, knowledge, and
moral weight. The structural identity is not a contingent coincidence but
a mathematical necessity: any symmetric bilinear form on which power,
knowledge, and ethics are defined must satisfy these identities.
\end{proof}

\begin{interpretation}[The Thesis of This Volume]
The structural identity theorem is the culmination of the five-chapter arc.
We have shown:
\begin{itemize}
\item In Chapter~11 (Kant), that the categorical imperative demands zero
emergence with the self---moral law as fixed point.
\item In Chapter~12 (Utilitarianism), that maximizing total moral weight
requires maximizing total self-resonance, but this program is undermined by
the non-commutativity of composition.
\item In Chapter~13 (Virtue Ethics), that character is accumulated moral
weight, built through iterated composition.
\item In Chapter~14 (Rawls and Justice), that the original position is a
resonance profile from which power has been subtracted.
\item In this chapter, that attribution of emergence is impossible, that
moral luck is sensitivity to context, that moral dilemmas compound rather
than resolve, and that forgiveness is creative emergence overcoming
negative resonance.
\end{itemize}

All of these results follow from the same axioms---the axioms of
$\mathrm{IdeaticSpace}$ developed in Volume I. The unity of power,
knowledge, and ethics is not a philosophical thesis that must be argued
for; it is a mathematical theorem that has been proved. The challenge
for human societies is not to \emph{discover} this unity but to build
institutions that \emph{respect} it: institutions that acknowledge that
every exercise of power is also an ethical act, that every act of
knowledge is also an exercise of power, and that the emergence created
by human composition belongs to no one and to everyone.
\end{interpretation}

\section{Conclusion: Power, Knowledge, and Ethics United}
\label{sec:ethics-conclusion}

\begin{remark}[The Unity of the Three]
The mathematics has spoken with unusual clarity. Power, knowledge, and
ethics are not three separate domains that happen to interact; they are
three perspectives on a single mathematical structure---the symmetric
bilinear form of resonance and the nonlinear operation of composition.

This unity has practical consequences:
\begin{enumerate}
\item \textbf{No ethics without power}: Moral weight is self-resonance,
and self-resonance is power. An agent with no power has no moral weight.
This is why disempowerment is not merely an injustice but a \emph{moral}
harm---it reduces the agent's capacity for moral standing.
\item \textbf{No power without ethics}: Every exercise of power creates
emergence, and emergence carries moral weight (by the enrichment axiom).
There is no morally neutral exercise of power. Even ``pure'' power---coercion,
domination, exclusion---creates moral weight that persists in the ideatic space.
\item \textbf{No knowledge without power or ethics}: Knowledge is mutual
resonance ($\rs(a,b) = \rs(b,a)$), and mutual resonance is a component of
both power and moral weight. To know something is simultaneously to be
empowered by it and to be morally implicated in it.
\end{enumerate}
\end{remark}

\begin{remark}[The Challenge Ahead]
The theorems proved in this volume are \emph{structural}---they describe
the mathematical constraints that any ethical theory must satisfy. They do
not tell us what to do; they tell us what the space of ethical possibilities
looks like. Within this space, we must still make choices: choices about
attribution, about justice, about forgiveness, about the tolerance of
intolerance.

But these choices are now \emph{informed} by mathematical structure. We know
that attribution is impossible (Theorem~\ref{thm:attribution-impossibility}),
that moral dilemmas compound (Theorem~\ref{thm:dilemma-compounds}), that
forgiveness requires creative emergence
(Theorem~\ref{thm:forgiveness-decomp}), and that power and moral weight are
structurally identical (Theorem~\ref{thm:structural-identity}).

The task of Volume VI (\emph{Emergence}) is to study the dynamics of these
structures---how they evolve over time, how emergence accumulates, how moral
weight grows. The task of the human reader is to decide what to do with
this knowledge.
\end{remark}
